\documentclass[../main.tex]{subfiles}
\begin{document}

\chapter{GPUアーキテクチャの最適化 I}
\lessoninfo{
  video={M6，レッスン 1--3（Handling Divergent Branch；Register Optimizations；Unified Register File and Shared Memory）},
  sources={Fung ら \cite{fung2007dwf}（動的ワープ形成）；Narasiman ら \cite{narasiman2011twolevel}（ラージワープ）；Gebhart ら \cite{gebhart2011energy}（レジスタファイルキャッシュ），\cite{gebhart2011compile}（コンパイラ管理のレジスタ階層），\cite{gebhart2012unifying}（統合ローカルメモリ）；Jeon ら \cite{jeon2015rfv}（レジスタファイル仮想化）；NVIDIA A100 \cite{nvidia2020ampere} および H100 \cite{nvidia2022hopper} ホワイトペーパー；Harris \cite{harris2013sharedmem}（共有メモリのレイテンシ）},
  textbook={A18 \cite{aamodt2018gpu} ch.~3（指定），ch.~4；K12 \cite{kim2012gpgpu} §1.4.1 GPU Pipeline，§1.4.2 Handling Branch Instructions；CPG \cite{nvidia2024cudaguide} §4.1 SIMT Architecture，Compute Capability 7.x--9.0 Shared Memory；PMPP \cite{hwu2022pmpp} ch.~10},
  keywords={ダイバージェント分岐，ダイバージェントワープ，再収束点，直接後支配節点，if変換，SIMT スタック，ラージワープ，動的ワープ形成，SIMD 効率，生存レジスタと非生存レジスタ，レジスタファイルキャッシュ，レジスタファイル仮想化，ソフトウェア管理キャッシュ，統合ローカルメモリ，共有メモリのカーブアウト},
}

第4章では SM のパイプラインを通るワープ命令をたどり，本章の土台となる
2つの事実を述べた．ワープはレーンを無効にできるアクティブマスクを持つこと，
そしてそのレジスタは，バンクに分割しなければならないほど大きな
レジスタファイルに置かれることである．本章の前半では，1つのワープの
スレッドが分岐の向きについて食い違ったときに何が起こるかを問う．
ハードウェアはそれをどう検出し，どのレーンがどの経路を実行し，
マスクとプログラムカウンタのスタックによって経路をどう再び合流させるか
である．続いて，ワープの構成を固定しないことでレーンの利用率を高める
2つの研究上のアイデアを述べる．次にレジスタファイルそのものを扱う．
それがなぜ大きく，遅く，電力を消費するのか，そしてレジスタの値に関する
4つの観察から何が言えるかを調べ，その手段として，前段に置く
レジスタファイルキャッシュと，名前の数より小さい仮想化されたファイルを
示す．最後にレジスタを共有メモリやキャッシュと比較し，1つのオンチップの
プールで3つすべてを受け持てるかを問う．これは NVIDIA GPU の構成可能な
L1 と共有メモリを研究として一般化したものである．第7章と同様に，本章は
研究論文を軸に構成されており，そうした論文を読むことが第2の目的である．

\section{ダイバージェント分岐の扱い}
\label{sec:l06-divergence}

第1章の SPMD モデルはすべてのスレッドで1つのプログラムを実行するが，
スレッドがその中を同じ経路でたどることを強制するものは何もない．\source{M6，
ダイバージェント分岐；K12 §1.4.2 Handling Branch Instructions；A18 ch.~3；
CPG §4.1 SIMT Architecture；\cite{fung2007dwf}．}本節ではダイバージェンスを
定義し，ダイバージェントな分岐を実行する2つの方法（プレディケーションと
SIMT スタック）を示し，そのコストに対する2つのマイクロアーキテクチャ上の
対策，すなわちラージワープと動的ワープ形成を紹介する．

\subsection{ダイバージェント分岐とワープ}

次のカーネルでは，各ブロックの最初の2スレッドがある処理を行い，残りが
別の処理を行う．本章の図のように4スレッドのワープを考えると，スレッド
0--3のワープは半分に分かれる．\margin{本章の図ではマスクを短く保つために
4スレッドのワープを描く．実際のワープサイズは，Compute Capability 5.0
から9.0までのすべての NVIDIA GPU で32である（CPG Table~21
\cite{nvidia2024cudaguide}）．}

\begin{lstlisting}[language=C++, numbers=none]
__global__ void split(float *a) {
  int i = threadIdx.x;      // BB1: 条件を計算する
  if (i < 2)
    a[i] = a[i] * 2.0f;     // BB2: Work-1
  else
    a[i] = a[i] + 1.0f;     // BB3: Work-2
  a[i] = a[i] - 0.5f;       // BB4: 再び両方の経路
}
\end{lstlisting}

コンパイラはこのようなコードを
\term[きほんぶろっく]{基本ブロック}[basic block]に切り分ける．
基本ブロックとは，最初の命令からしか入らず，最後の命令の後でしか
出ない極大の命令列であり，その命令はすべて実行されるか，どれも
実行されないかのどちらかである．ブロックと，それらの間の可能な制御の
移動が\cref{fig:l06-divergent-cfg}の制御フローグラフを成す．
制御フローグラフは第10章で形式的に定義する．

\begin{definition}[ダイバージェント分岐とダイバージェントワープ]
  \label{def:l06-divergent}
  \term[だいばーじぇんとぶんき]{ダイバージェント分岐}[divergent branch]
  とは，その結果が1つのワープのスレッドの間で異なり，その後ワープの
  スレッドが異なるプログラムカウンタからフェッチすることになる分岐である．
  このときそのワープは
  \term[だいばーじぇんとわーぷ]{ダイバージェントワープ}[divergent warp]
  であり，この現象を
  \term[わーぷだいばーじぇんす]{ワープダイバージェンス}[warp divergence]
  という．すべてのワープのすべてのスレッドが一致する分岐はワープ一様で
  あり，スカラプロセッサ上と同じコストしかかからない．
\end{definition}

\begin{figure}[tb]
  \centering
  % A divergent branch (Lesson 6).  Left: the control-flow graph of
% "if (i < 2) Work-1 else Work-2" for a four-thread warp, each basic
% block with its active mask; the paths split after BB1 and reconverge
% at BB4.  Right: the lanes that execute each block, in time order.
% Redraws the outline figure m6l1_divb (Work-2 -> BB4 edge added).
% Usage: % A divergent branch (Lesson 6).  Left: the control-flow graph of
% "if (i < 2) Work-1 else Work-2" for a four-thread warp, each basic
% block with its active mask; the paths split after BB1 and reconverge
% at BB4.  Right: the lanes that execute each block, in time order.
% Redraws the outline figure m6l1_divb (Work-2 -> BB4 edge added).
% Usage: % A divergent branch (Lesson 6).  Left: the control-flow graph of
% "if (i < 2) Work-1 else Work-2" for a four-thread warp, each basic
% block with its active mask; the paths split after BB1 and reconverge
% at BB4.  Right: the lanes that execute each block, in time order.
% Redraws the outline figure m6l1_divb (Work-2 -> BB4 edge added).
% Usage: \input{figures/l06-divergent-cfg}
% Width: body (120 mm); the drawing is about 108 mm wide.
\begin{tikzpicture}[x=1mm, y=1mm, font=\scriptsize\sffamily,
    >={Stealth[length=1.6mm]},
    bb/.style={draw, minimum width=22mm, minimum height=8mm,
               align=center, inner sep=1pt}]
  % CFG
  \node[bb, fill=ln-accent!20] (b1) at (24,42)
    {BB1: \texttt{if (i < 2)}\\\texttt{1111}};
  \node[bb, fill=ln-source!25] (b2) at (11,26)
    {BB2: Work-1\\\texttt{1100}};
  \node[bb, fill=ln-tip!25] (b3) at (37,26)
    {BB3: Work-2\\\texttt{0011}};
  \node[bb, fill=ln-accent!20] (b4) at (24,10)
    {BB4\\\texttt{1111}};
  \draw[->] (b1) -- node[left, inner sep=2pt] {\iflangja{真}{taken}} (b2);
  \draw[->] (b1) -- node[right, inner sep=2pt] {\iflangja{偽}{not taken}} (b3);
  \draw[->] (b2) -- (b4);
  \draw[->] (b3) -- (b4);
  \node[anchor=north, text=ln-muted, align=center] at (24,4.5)
    {\iflangja{BB4 = 再収束点}{BB4 = reconvergence point}};
  % lanes over time
  \node[anchor=south west, font=\scriptsize\sffamily\bfseries] at (58,47)
    {\iflangja{実行するレーン}{lanes that execute}};
  \foreach \l in {0,...,3}
    \node[text=ln-muted] at (73+\l*7,44) {\iflangja{\l}{lane \l}};
  \foreach \row/\name/\m/\c in {0/BB1/{1,1,1,1}/ln-accent!20, 1/BB2/{1,1,0,0}/ln-source!25,
                                 2/BB3/{0,0,1,1}/ln-tip!25, 3/BB4/{1,1,1,1}/ln-accent!20} {
    \node[anchor=east] at (69,38-\row*9) {\name};
    \foreach \bit [count=\l from 0] in \m {
      \ifnum\bit=1
        \draw[fill=\c] (70.5+\l*7,35.5-\row*9) rectangle ++(5,5);
      \else
        \draw[ln-rule, dashed] (70.5+\l*7,35.5-\row*9) rectangle ++(5,5);
      \fi
    }
  }
  \draw[->, ln-muted] (101,41) -- (101,9) node[midway, right, inner sep=1pt]
    {\iflangja{時間}{time}};
\end{tikzpicture}

% Width: body (120 mm); the drawing is about 108 mm wide.
\begin{tikzpicture}[x=1mm, y=1mm, font=\scriptsize\sffamily,
    >={Stealth[length=1.6mm]},
    bb/.style={draw, minimum width=22mm, minimum height=8mm,
               align=center, inner sep=1pt}]
  % CFG
  \node[bb, fill=ln-accent!20] (b1) at (24,42)
    {BB1: \texttt{if (i < 2)}\\\texttt{1111}};
  \node[bb, fill=ln-source!25] (b2) at (11,26)
    {BB2: Work-1\\\texttt{1100}};
  \node[bb, fill=ln-tip!25] (b3) at (37,26)
    {BB3: Work-2\\\texttt{0011}};
  \node[bb, fill=ln-accent!20] (b4) at (24,10)
    {BB4\\\texttt{1111}};
  \draw[->] (b1) -- node[left, inner sep=2pt] {\iflangja{真}{taken}} (b2);
  \draw[->] (b1) -- node[right, inner sep=2pt] {\iflangja{偽}{not taken}} (b3);
  \draw[->] (b2) -- (b4);
  \draw[->] (b3) -- (b4);
  \node[anchor=north, text=ln-muted, align=center] at (24,4.5)
    {\iflangja{BB4 = 再収束点}{BB4 = reconvergence point}};
  % lanes over time
  \node[anchor=south west, font=\scriptsize\sffamily\bfseries] at (58,47)
    {\iflangja{実行するレーン}{lanes that execute}};
  \foreach \l in {0,...,3}
    \node[text=ln-muted] at (73+\l*7,44) {\iflangja{\l}{lane \l}};
  \foreach \row/\name/\m/\c in {0/BB1/{1,1,1,1}/ln-accent!20, 1/BB2/{1,1,0,0}/ln-source!25,
                                 2/BB3/{0,0,1,1}/ln-tip!25, 3/BB4/{1,1,1,1}/ln-accent!20} {
    \node[anchor=east] at (69,38-\row*9) {\name};
    \foreach \bit [count=\l from 0] in \m {
      \ifnum\bit=1
        \draw[fill=\c] (70.5+\l*7,35.5-\row*9) rectangle ++(5,5);
      \else
        \draw[ln-rule, dashed] (70.5+\l*7,35.5-\row*9) rectangle ++(5,5);
      \fi
    }
  }
  \draw[->, ln-muted] (101,41) -- (101,9) node[midway, right, inner sep=1pt]
    {\iflangja{時間}{time}};
\end{tikzpicture}

% Width: body (120 mm); the drawing is about 108 mm wide.
\begin{tikzpicture}[x=1mm, y=1mm, font=\scriptsize\sffamily,
    >={Stealth[length=1.6mm]},
    bb/.style={draw, minimum width=22mm, minimum height=8mm,
               align=center, inner sep=1pt}]
  % CFG
  \node[bb, fill=ln-accent!20] (b1) at (24,42)
    {BB1: \texttt{if (i < 2)}\\\texttt{1111}};
  \node[bb, fill=ln-source!25] (b2) at (11,26)
    {BB2: Work-1\\\texttt{1100}};
  \node[bb, fill=ln-tip!25] (b3) at (37,26)
    {BB3: Work-2\\\texttt{0011}};
  \node[bb, fill=ln-accent!20] (b4) at (24,10)
    {BB4\\\texttt{1111}};
  \draw[->] (b1) -- node[left, inner sep=2pt] {\iflangja{真}{taken}} (b2);
  \draw[->] (b1) -- node[right, inner sep=2pt] {\iflangja{偽}{not taken}} (b3);
  \draw[->] (b2) -- (b4);
  \draw[->] (b3) -- (b4);
  \node[anchor=north, text=ln-muted, align=center] at (24,4.5)
    {\iflangja{BB4 = 再収束点}{BB4 = reconvergence point}};
  % lanes over time
  \node[anchor=south west, font=\scriptsize\sffamily\bfseries] at (58,47)
    {\iflangja{実行するレーン}{lanes that execute}};
  \foreach \l in {0,...,3}
    \node[text=ln-muted] at (73+\l*7,44) {\iflangja{\l}{lane \l}};
  \foreach \row/\name/\m/\c in {0/BB1/{1,1,1,1}/ln-accent!20, 1/BB2/{1,1,0,0}/ln-source!25,
                                 2/BB3/{0,0,1,1}/ln-tip!25, 3/BB4/{1,1,1,1}/ln-accent!20} {
    \node[anchor=east] at (69,38-\row*9) {\name};
    \foreach \bit [count=\l from 0] in \m {
      \ifnum\bit=1
        \draw[fill=\c] (70.5+\l*7,35.5-\row*9) rectangle ++(5,5);
      \else
        \draw[ln-rule, dashed] (70.5+\l*7,35.5-\row*9) rectangle ++(5,5);
      \fi
    }
  }
  \draw[->, ln-muted] (101,41) -- (101,9) node[midway, right, inner sep=1pt]
    {\iflangja{時間}{time}};
\end{tikzpicture}

  \caption{4スレッドのワープにおけるダイバージェント分岐．左：
    制御フローグラフ．各基本ブロックに，それを実行するときの
    アクティブマスクを添える（左から $k$ 番目のビットがレーン~$k$）．
    右：各ブロックを実行するレーン．BB2 と BB3 は順に実行され，
    それぞれレーンの半分が遊ぶ．}
  \label{fig:l06-divergent-cfg}
\end{figure}

第4章のアクティブマスクは，各ブロックにどのレーンが参加するかを記録する
（\cref{fig:l06-divergent-cfg}）．BB1 は \texttt{1111}，BB2 は
\texttt{1100}，BB3 は \texttt{0011}，BB4 は再び \texttt{1111} で実行される．
ワープのプログラムカウンタは1つなので，2つの経路を同時には実行できない．
それらは順に実行され，そのたびにもう一方の経路のレーンは無効にされ，
2つの経路の時間は足し合わされる．\caution{ダイバージェント分岐は
コードを誤らせるのではなく，遅くするだけである．結果はスカラ機械が
計算するものと同じである．\emph{異なる}ワープの間のダイバージェンスには
コストがかからない．各ワープが自分のプログラムカウンタを持つからである．}
経路が再び出会う点 BB4 で，ワープは完全なマスクで再開できる．

\begin{definition}[再収束点]
  \label{def:l06-reconvergence}
  ダイバージェント分岐の\term[さいしゅうそくてん]{再収束点}[reconvergence point]
  とは，分岐から出たワープのすべての経路が再び出会うことが保証される
  最初の命令であり，そこからワープは全体として進むことができる．
  コンパイラの用語では，ブロック $X$ からプログラムの出口へ至るすべての
  経路がブロック $Y$ を通るとき，$Y$ は $X$ の
  \term[こうしはいせってん]{後支配節点}[post-dominator]であり，
  $X$ の
  \term[ちょくせつこうしはいせってん]{直接後支配節点}[immediate post-dominator (IPDOM)]\index{IPDOM|see{直接後支配節点}}（IPDOM）
  は $X$ 自身を除いて最も近い後支配節点である．分岐ブロックの IPDOM が
  通常の再収束点であり，\cref{fig:l06-divergent-cfg}ではそれは BB4 である．
\end{definition}

\margin{Fung らは，ワープの分かれた部分をカーネルの終わりまで別々に
走らせるのではなく直接後支配節点で再収束させると，93.4\,\% の高速化が
得られると報告している \cite{fung2007dwf}．}%
できるだけ早く再収束することが重要なのは，分岐の後，再収束点の前に
実行されるすべてのブロックが部分的なマスクで実行されるからである．

\subsection{if変換とプレディケーション}

分岐を取り除く最も古い方法は，ジャンプを取り除くことである．

\begin{definition}[if変換]
  \label{def:l06-if-conversion}
  \term[いふへんかん]{if変換}[if-conversion]と呼ばれる変換は，分岐と
  その2つの経路を直線的なコードで置き換える．1つの命令が条件を
  \term[ぷれでぃけーと]{プレディケート}[predicate]，すなわちレーンごとの
  1ビットのレジスタに計算し，元の経路の各命令はそのプレディケートまたは
  その否定によってガードされる（\cref{fig:l06-if-conversion}）．
  ガードが偽のレーンは結果を書き込まない．得られたコードは第5章の
  プレディケート実行によって実行される．
\end{definition}

\begin{marginfigure}
  \resizebox{\linewidth}{!}{% If-conversion (Lesson 6): the branch of figure l06-divergent-cfg
% becomes one straight-line block of predicated instructions that every
% lane fetches; no lane ever jumps.
% Redraws the outline figure m6l1_divb2.
% Usage: % If-conversion (Lesson 6): the branch of figure l06-divergent-cfg
% becomes one straight-line block of predicated instructions that every
% lane fetches; no lane ever jumps.
% Redraws the outline figure m6l1_divb2.
% Usage: % If-conversion (Lesson 6): the branch of figure l06-divergent-cfg
% becomes one straight-line block of predicated instructions that every
% lane fetches; no lane ever jumps.
% Redraws the outline figure m6l1_divb2.
% Usage: \input{figures/l06-if-conversion}  (in a marginfigure)
% Width: margin (50 mm); the drawing is about 46 mm wide.
\begin{tikzpicture}[x=1mm, y=1mm, font=\scriptsize\sffamily,
    >={Stealth[length=1.6mm]},
    ins/.style={draw, minimum width=34mm, minimum height=6mm,
                align=center, inner sep=1pt}]
  \node[ins, fill=ln-accent!20] (p) at (0,30)
    {\iflangja{プレディケート\\\texttt{p = (i < 2)}}{set predicate \texttt{p = (i < 2)}}};
  \node[ins, fill=ln-source!25] (a) at (0,21) {\texttt{(p)} Work-1};
  \node[ins, fill=ln-tip!25] (b) at (0,12) {\texttt{(!p)} Work-2};
  \node[ins, fill=ln-accent!20] (c) at (0,3) {BB4};
  \draw[->] (p) -- (a);
  \draw[->] (a) -- (b);
  \draw[->] (b) -- (c);
  \draw[ln-muted] (19,33) -- ++(1.5,0) -- ++(0,-33) -- ++(-1.5,0);
  \node[rotate=90, text=ln-muted] at (23.5,16.5)
    {\iflangja{全レーンがフェッチ}{fetched by every lane}};
\end{tikzpicture}
  (in a marginfigure)
% Width: margin (50 mm); the drawing is about 46 mm wide.
\begin{tikzpicture}[x=1mm, y=1mm, font=\scriptsize\sffamily,
    >={Stealth[length=1.6mm]},
    ins/.style={draw, minimum width=34mm, minimum height=6mm,
                align=center, inner sep=1pt}]
  \node[ins, fill=ln-accent!20] (p) at (0,30)
    {\iflangja{プレディケート\\\texttt{p = (i < 2)}}{set predicate \texttt{p = (i < 2)}}};
  \node[ins, fill=ln-source!25] (a) at (0,21) {\texttt{(p)} Work-1};
  \node[ins, fill=ln-tip!25] (b) at (0,12) {\texttt{(!p)} Work-2};
  \node[ins, fill=ln-accent!20] (c) at (0,3) {BB4};
  \draw[->] (p) -- (a);
  \draw[->] (a) -- (b);
  \draw[->] (b) -- (c);
  \draw[ln-muted] (19,33) -- ++(1.5,0) -- ++(0,-33) -- ++(-1.5,0);
  \node[rotate=90, text=ln-muted] at (23.5,16.5)
    {\iflangja{全レーンがフェッチ}{fetched by every lane}};
\end{tikzpicture}
  (in a marginfigure)
% Width: margin (50 mm); the drawing is about 46 mm wide.
\begin{tikzpicture}[x=1mm, y=1mm, font=\scriptsize\sffamily,
    >={Stealth[length=1.6mm]},
    ins/.style={draw, minimum width=34mm, minimum height=6mm,
                align=center, inner sep=1pt}]
  \node[ins, fill=ln-accent!20] (p) at (0,30)
    {\iflangja{プレディケート\\\texttt{p = (i < 2)}}{set predicate \texttt{p = (i < 2)}}};
  \node[ins, fill=ln-source!25] (a) at (0,21) {\texttt{(p)} Work-1};
  \node[ins, fill=ln-tip!25] (b) at (0,12) {\texttt{(!p)} Work-2};
  \node[ins, fill=ln-accent!20] (c) at (0,3) {BB4};
  \draw[->] (p) -- (a);
  \draw[->] (a) -- (b);
  \draw[->] (b) -- (c);
  \draw[ln-muted] (19,33) -- ++(1.5,0) -- ++(0,-33) -- ++(-1.5,0);
  \node[rotate=90, text=ln-muted] at (23.5,16.5)
    {\iflangja{全レーンがフェッチ}{fetched by every lane}};
\end{tikzpicture}
}
  \caption{if変換後の\cref{fig:l06-divergent-cfg}の分岐．
    ガード付きの命令を全レーンがフェッチする1つの基本ブロック．}
  \label{fig:l06-if-conversion}
\end{marginfigure}

if変換はベクトルプロセッサでは標準的であり，CPU も短い分岐本体に
用いる．そこでは条件付き移動のほうが，起こりうる分岐予測ミスより安い．
GPU では if変換はダイバージェンスを完全に取り除く．どのレーンも
ジャンプしないので，ワープは1つのプログラムカウンタを保つ．その代償は，
どちらかの経路をどのレーンも必要としないときでも，両方の経路の命令が
フェッチされ発行されることである．そのためコンパイラは if変換を短い
本体にしか適用しない．\margin{PTX ではガードは命令の前に \texttt{@p}
または \texttt{@!p} と書く．プレディケートレジスタは第10章で示す．}

\begin{example}
  \label{ex:l06-if-conversion}
  上のカーネルで，BB1 と BB4 をそれぞれ2命令，BB2 と BB3 をそれぞれ
  6命令とする．4レーンのワープを考え，発行された命令1つにつきレーン1つ
  あたり1つとしてレーンスロットを数える．
  \begin{itemize}
    \item \emph{ダイバージェントなワープ，分岐コード}：発行命令は
      $2+6+6+2 = 16$，レーンスロットは $64$ で，そのうち $2\cdot4 + 6\cdot2 +
      6\cdot2 + 2\cdot4 = 40$ が仕事をする．
    \item \emph{ダイバージェントなワープ，if変換後}：同じ16命令
      （プレディケートは BB1 で計算する）と，40の有用なレーンスロット．
    \item \emph{一様なワープ（全レーンが BB2 を取る），分岐コード}：
      $2+6+2 = 10$ 命令で，全レーンが働く．
    \item \emph{一様なワープ，if変換後}：やはり16命令である．BB3 の
      6命令はどのレーンでも何もしない．
  \end{itemize}
  if変換はコストをデータに依存しないものにする．この計数では，ワープが
  ダイバージェントなら引き分けで，ワープが一様なら負ける．if変換が
  節約するのはこの計数が除外したオーバーヘッド（分岐命令，スタックの
  プッシュとポップ，再収束の検査，フェッチ先の変更）であり，それが
  効くのは短い本体の場合だけである．\tip{if変換のコスト：常に両方の
本体の和．分岐のコスト：ワープが一様なら取った側の本体，
ダイバージェントなら両方の和．}
\end{example}

\subsection{SIMT スタック}
\label{sec:l06-simt-stack}

プレディケーションを使わないなら，ハードウェア自身が両方の経路を実行し，
それらを合流させなければならない．\caution{講義 vs 教科書：以下で組み立てる
2フィールドのエントリでは，再収束点で完全なマスクがどう戻るかが
暗黙のままである．K12 §1.4.2 はエントリに3つのフィールド，
\enquote{the PC address of the reconvergence point, active mask bits,
and the next PC address} を与えている \cite{kim2012gpgpu}．}
この機構は，\cref{fig:l06-divergent-cfg}のワープをたどりながら段階的に
組み立てると最も理解しやすい．
\begin{enumerate}
  \item BB1 の後，4つのスレッドの次の PC は異なる（レーン 0--1 は BB2，
    レーン 2--3 は BB3）．分岐はダイバージェントである．1つの PC と
    1つのマスクでは両方の行き先を表せない．
  \item 1つの経路（BB1，BB2，BB4）だけをたどるのは誤りである．
    レーン 2--3 のスレッドが BB3 を実行しないからである．
  \item 解決策は取らなかった経路を覚えておくことである．すなわち
    \emph{代替 PC}（BB3）とそのマスク（\texttt{0011}）を，取った経路の
    後に実行する．順序は BB1，BB2，BB3，BB4 となる．
  \item 分岐が入れ子になると，代替スロット1つでは足りない．BB2 自体が
    ダイバージェントになると（\cref{fig:l06-simt-stack}左：BB4 は
    \texttt{1000}，BB5 は \texttt{0100}），その代替 PC BB5 が唯一のスロットの
    BB3 を上書きし，BB3 は失われる．
  \item そこで保留中の経路をスタックに積む．
\end{enumerate}

\begin{definition}[SIMT スタック]
  \label{def:l06-simt-stack}
  ワープの\term[えすあいえむてぃーすたっく]{SIMTスタック}[SIMT stack]は，
  ダイバージェント分岐の保留中の経路ごとに1つのエントリを保持する．
  エントリは代替 PC とそのアクティブマスク，そして現在の経路が終わる
  再収束 PC である．ダイバージェント分岐でハードウェアは取らなかった
  経路のエントリをプッシュし，取った経路を縮小したマスクで続ける．
  フェッチのたびに，次の PC を先頭エントリの再収束 PC と比較し，等しければ
  そのエントリを\emph{ポップ}して，代わりにエントリのマスクで代替 PC から
  フェッチする．\term[いれこのぶんき]{入れ子の分岐}[nested branch]，
  すなわち別の分岐の経路の中にあるダイバージェント分岐は，単にもう1つ
  エントリをプッシュする．
\end{definition}

\begin{figure}[tb]
  \centering
  % The SIMT stack on a nested branch (Lesson 6).  Left: a CFG in which
% BB2 branches again (BB4 1000, BB5 0100) and all paths reconverge at
% BB6.  Right: the stack of (alternate PC with its mask, reconvergence
% PC) entries after BB1, after BB2, after the pop at the end of BB4 and
% after the pop at the end of BB5 (top entry first).
% Redraws the outline figures m6l1_gpuflow4 and m6l1_gpuflow5 (BB2 mask
% corrected to 1100).
% Usage: % The SIMT stack on a nested branch (Lesson 6).  Left: a CFG in which
% BB2 branches again (BB4 1000, BB5 0100) and all paths reconverge at
% BB6.  Right: the stack of (alternate PC with its mask, reconvergence
% PC) entries after BB1, after BB2, after the pop at the end of BB4 and
% after the pop at the end of BB5 (top entry first).
% Redraws the outline figures m6l1_gpuflow4 and m6l1_gpuflow5 (BB2 mask
% corrected to 1100).
% Usage: % The SIMT stack on a nested branch (Lesson 6).  Left: a CFG in which
% BB2 branches again (BB4 1000, BB5 0100) and all paths reconverge at
% BB6.  Right: the stack of (alternate PC with its mask, reconvergence
% PC) entries after BB1, after BB2, after the pop at the end of BB4 and
% after the pop at the end of BB5 (top entry first).
% Redraws the outline figures m6l1_gpuflow4 and m6l1_gpuflow5 (BB2 mask
% corrected to 1100).
% Usage: \input{figures/l06-simt-stack}
% Width: body (120 mm); the drawing is about 112 mm wide.
\begin{tikzpicture}[x=1mm, y=1mm, font=\scriptsize\sffamily,
    >={Stealth[length=1.6mm]},
    bb/.style={draw, minimum width=13mm, minimum height=7.5mm,
               align=center, inner sep=1pt}]
  % nested CFG
  \node[bb, fill=ln-accent!20] (b1) at (22,56) {BB1\\\texttt{1111}};
  \node[bb, fill=ln-source!25] (b2) at (12,41) {BB2\\\texttt{1100}};
  \node[bb, fill=ln-tip!25] (b3) at (36,41) {BB3\\\texttt{0011}};
  \node[bb, fill=ln-source!12] (b4) at (4,25) {BB4\\\texttt{1000}};
  \node[bb, fill=ln-source!40] (b5) at (20,25) {BB5\\\texttt{0100}};
  \node[bb, fill=ln-accent!20] (b6) at (22,9) {BB6\\\texttt{1111}};
  \draw[->] (b1) -- (b2);
  \draw[->] (b1) -- (b3);
  \draw[->] (b2) -- (b4);
  \draw[->] (b2) -- (b5);
  \draw[->] (b4) -- (b6);
  \draw[->] (b5) -- (b6);
  \draw[->] (b3) |- (b6);
  % stack snapshots: \snap{x}{y}{title}{list of alt/reconv}
  \newcommand{\snapshot}[4]{%
    \node[anchor=south west, font=\scriptsize\sffamily\bfseries] at (#1,#2+1) {#3};
    \node[anchor=south, text=ln-muted] at (#1+8,#2-4.5) {\iflangja{代替 PC}{alt.\ PC}};
    \node[anchor=south, text=ln-muted] at (#1+21.5,#2-4.5) {\iflangja{再収束点}{reconv.}};
    \foreach \alt/\rec [count=\k from 0] in {#4} {
      \draw[fill=ln-box] (#1,#2-5-\k*6) rectangle ++(16,-5.5)
        node[midway] {\alt};
      \draw[fill=ln-box] (#1+16,#2-5-\k*6) rectangle ++(11,-5.5)
        node[midway] {\rec};
    }}
  \snapshot{52}{56}{\iflangja{(a) BB1 の後}{(a) after BB1}}
    {{BB3 \texttt{0011}}/BB6}
  \snapshot{84}{56}{\iflangja{(b) BB2 の後}{(b) after BB2}}
    {{BB5 \texttt{0100}}/BB6, {BB3 \texttt{0011}}/BB6}
  \snapshot{52}{26}{\iflangja{(c) BB4 の後（pop）}{(c) after BB4 (pop)}}
    {{BB3 \texttt{0011}}/BB6}
  \node[anchor=south west, font=\scriptsize\sffamily\bfseries] at (84,27)
    {\iflangja{(d) BB5 の後（pop）}{(d) after BB5 (pop)}};
  \node[anchor=north west, text=ln-muted, align=left] at (84,24)
    {\iflangja{空：BB3 の後は\\BB6 を \texttt{1111} で実行}{empty: after BB3,\\BB6 runs with \texttt{1111}}};
\end{tikzpicture}

% Width: body (120 mm); the drawing is about 112 mm wide.
\begin{tikzpicture}[x=1mm, y=1mm, font=\scriptsize\sffamily,
    >={Stealth[length=1.6mm]},
    bb/.style={draw, minimum width=13mm, minimum height=7.5mm,
               align=center, inner sep=1pt}]
  % nested CFG
  \node[bb, fill=ln-accent!20] (b1) at (22,56) {BB1\\\texttt{1111}};
  \node[bb, fill=ln-source!25] (b2) at (12,41) {BB2\\\texttt{1100}};
  \node[bb, fill=ln-tip!25] (b3) at (36,41) {BB3\\\texttt{0011}};
  \node[bb, fill=ln-source!12] (b4) at (4,25) {BB4\\\texttt{1000}};
  \node[bb, fill=ln-source!40] (b5) at (20,25) {BB5\\\texttt{0100}};
  \node[bb, fill=ln-accent!20] (b6) at (22,9) {BB6\\\texttt{1111}};
  \draw[->] (b1) -- (b2);
  \draw[->] (b1) -- (b3);
  \draw[->] (b2) -- (b4);
  \draw[->] (b2) -- (b5);
  \draw[->] (b4) -- (b6);
  \draw[->] (b5) -- (b6);
  \draw[->] (b3) |- (b6);
  % stack snapshots: \snap{x}{y}{title}{list of alt/reconv}
  \newcommand{\snapshot}[4]{%
    \node[anchor=south west, font=\scriptsize\sffamily\bfseries] at (#1,#2+1) {#3};
    \node[anchor=south, text=ln-muted] at (#1+8,#2-4.5) {\iflangja{代替 PC}{alt.\ PC}};
    \node[anchor=south, text=ln-muted] at (#1+21.5,#2-4.5) {\iflangja{再収束点}{reconv.}};
    \foreach \alt/\rec [count=\k from 0] in {#4} {
      \draw[fill=ln-box] (#1,#2-5-\k*6) rectangle ++(16,-5.5)
        node[midway] {\alt};
      \draw[fill=ln-box] (#1+16,#2-5-\k*6) rectangle ++(11,-5.5)
        node[midway] {\rec};
    }}
  \snapshot{52}{56}{\iflangja{(a) BB1 の後}{(a) after BB1}}
    {{BB3 \texttt{0011}}/BB6}
  \snapshot{84}{56}{\iflangja{(b) BB2 の後}{(b) after BB2}}
    {{BB5 \texttt{0100}}/BB6, {BB3 \texttt{0011}}/BB6}
  \snapshot{52}{26}{\iflangja{(c) BB4 の後（pop）}{(c) after BB4 (pop)}}
    {{BB3 \texttt{0011}}/BB6}
  \node[anchor=south west, font=\scriptsize\sffamily\bfseries] at (84,27)
    {\iflangja{(d) BB5 の後（pop）}{(d) after BB5 (pop)}};
  \node[anchor=north west, text=ln-muted, align=left] at (84,24)
    {\iflangja{空：BB3 の後は\\BB6 を \texttt{1111} で実行}{empty: after BB3,\\BB6 runs with \texttt{1111}}};
\end{tikzpicture}

% Width: body (120 mm); the drawing is about 112 mm wide.
\begin{tikzpicture}[x=1mm, y=1mm, font=\scriptsize\sffamily,
    >={Stealth[length=1.6mm]},
    bb/.style={draw, minimum width=13mm, minimum height=7.5mm,
               align=center, inner sep=1pt}]
  % nested CFG
  \node[bb, fill=ln-accent!20] (b1) at (22,56) {BB1\\\texttt{1111}};
  \node[bb, fill=ln-source!25] (b2) at (12,41) {BB2\\\texttt{1100}};
  \node[bb, fill=ln-tip!25] (b3) at (36,41) {BB3\\\texttt{0011}};
  \node[bb, fill=ln-source!12] (b4) at (4,25) {BB4\\\texttt{1000}};
  \node[bb, fill=ln-source!40] (b5) at (20,25) {BB5\\\texttt{0100}};
  \node[bb, fill=ln-accent!20] (b6) at (22,9) {BB6\\\texttt{1111}};
  \draw[->] (b1) -- (b2);
  \draw[->] (b1) -- (b3);
  \draw[->] (b2) -- (b4);
  \draw[->] (b2) -- (b5);
  \draw[->] (b4) -- (b6);
  \draw[->] (b5) -- (b6);
  \draw[->] (b3) |- (b6);
  % stack snapshots: \snap{x}{y}{title}{list of alt/reconv}
  \newcommand{\snapshot}[4]{%
    \node[anchor=south west, font=\scriptsize\sffamily\bfseries] at (#1,#2+1) {#3};
    \node[anchor=south, text=ln-muted] at (#1+8,#2-4.5) {\iflangja{代替 PC}{alt.\ PC}};
    \node[anchor=south, text=ln-muted] at (#1+21.5,#2-4.5) {\iflangja{再収束点}{reconv.}};
    \foreach \alt/\rec [count=\k from 0] in {#4} {
      \draw[fill=ln-box] (#1,#2-5-\k*6) rectangle ++(16,-5.5)
        node[midway] {\alt};
      \draw[fill=ln-box] (#1+16,#2-5-\k*6) rectangle ++(11,-5.5)
        node[midway] {\rec};
    }}
  \snapshot{52}{56}{\iflangja{(a) BB1 の後}{(a) after BB1}}
    {{BB3 \texttt{0011}}/BB6}
  \snapshot{84}{56}{\iflangja{(b) BB2 の後}{(b) after BB2}}
    {{BB5 \texttt{0100}}/BB6, {BB3 \texttt{0011}}/BB6}
  \snapshot{52}{26}{\iflangja{(c) BB4 の後（pop）}{(c) after BB4 (pop)}}
    {{BB3 \texttt{0011}}/BB6}
  \node[anchor=south west, font=\scriptsize\sffamily\bfseries] at (84,27)
    {\iflangja{(d) BB5 の後（pop）}{(d) after BB5 (pop)}};
  \node[anchor=north west, text=ln-muted, align=left] at (84,24)
    {\iflangja{空：BB3 の後は\\BB6 を \texttt{1111} で実行}{empty: after BB3,\\BB6 runs with \texttt{1111}}};
\end{tikzpicture}

  \caption{入れ子の分岐における SIMT スタック．左：BB2 が再び分岐し，
    すべての経路が BB6 で再収束する．右：BB1 の分岐の後，BB2 の分岐の後，
    BB4 の終わりのポップの後，BB5 の終わりのポップの後のスタック
    （先頭エントリが上）．}
  \label{fig:l06-simt-stack}
\end{figure}

\cref{tab:l06-simt-trace}は入れ子の例をたどる．BB4 は次の PC BB6 で
終わり，これはスタックの先頭の再収束 PC に等しいので，ワープはポップして
BB5 を \texttt{0100} でフェッチする．BB5 も BB6 で終わるので，ワープは
2番目のエントリをポップし，BB3 を \texttt{0011} でフェッチする．BB3 が
終わるとスタックは空になり，BB6 が完全なマスクで実行される．
2フィールドのスタックは \texttt{1111} がどこから来るかを示さない．
傍注の3フィールドのスタックは，分岐前のエントリからそれを復元する．
すべてのスレッドがちょうど自分の経路を実行し，6つのブロックに6回の
フェッチを要した．\sidenote{K12 §1.4.2 と A18 ch.~3
\cite{kim2012gpgpu,aamodt2018gpu} の3フィールドの定式化では，スタックの
先頭は常にワープが実行しているものを表す．ダイバージェント分岐では
先頭エントリの次 PC が再収束 PC に置き換えられ（分岐前のマスクは保つ），
経路ごとに1つのエントリが，それぞれ自分の次 PC，マスク，共通の再収束 PC
とともにプッシュされる．ワープが先頭エントリの再収束 PC に達すると
そのエントリがポップされ，次のエントリが PC とマスクを与える．
したがって BB6 では分岐前のエントリが \texttt{1111} を復元する．}

\begin{table}[tb]
  \centering
  \caption{\cref{fig:l06-simt-stack}の入れ子の分岐を SIMT スタックで
    実行する様子．スタックのエントリは（代替 PC，マスク；再収束 PC）と
    書き，先頭を最初に置く．}
  \label{tab:l06-simt-trace}
  \setlength{\tabcolsep}{4pt}
  \begin{tabular}{@{}l l l >{\raggedright\arraybackslash}p{34mm} >{\raggedright\arraybackslash}p{33mm}@{}}
    \toprule
    ブロック & マスク & 次の PC & 動作 & 実行後のスタック \\
    \midrule
    BB1 & \texttt{1111} & BB2 / BB3 & 分岐：BB3 をプッシュ，BB2 へ
      & (BB3, \texttt{0011}; BB6) \\
    BB2 & \texttt{1100} & BB4 / BB5 & 分岐：BB5 をプッシュ，BB4 へ
      & (BB5, \texttt{0100}; BB6), (BB3, \texttt{0011}; BB6) \\
    BB4 & \texttt{1000} & BB6 & 先頭の再収束点と一致：ポップし BB5 へ
      & (BB3, \texttt{0011}; BB6) \\
    BB5 & \texttt{0100} & BB6 & 先頭の再収束点と一致：ポップし BB3 へ
      & 空 \\
    BB3 & \texttt{0011} & BB6 & スタックが空：完全なマスク（暗黙） & 空 \\
    BB6 & \texttt{1111} & --- & --- & 空 \\
    \bottomrule
  \end{tabular}
\end{table}

スタックのハードウェアコストは，ワープあたり数個のエントリと，フェッチ
ごとの再収束の比較である．保持すべきエントリ数はダイバージェント分岐の
入れ子の深さとともに増える．

\subsection{再収束点の求め方}

再収束 PC はコンパイル時に計算される．コンパイラは制御フローグラフを
構築し，各分岐の直接後支配節点を求め，それをコードに印すか，スタックを
プッシュしポップする明示的な命令を挿入する．スタック自体はハードウェアに
置かれる．どの分岐が実際にダイバージェントになるかは実行時にしか
分からないからである．原理的にはコンパイラは決してダイバージェントに
ならない分岐についてスタック操作を省けるが，分岐がダイバージェントに
なるかを静的に判定するのは難しい．そのためコンパイラは，そうでないと
証明できない限り分岐をダイバージェントとして扱わねばならず，
ダイバージェント分岐を過大に報告する．\margin{第11章では
ダイバージェンス解析を扱う．これは，どの値が，したがってどの分岐条件が
ワープのスレッド間で異なりうるかを決めるコンパイラの解析である．}

\begin{note}
  SIMT スタックは分岐の経路を逐次化する．ワープは一方の経路を
  再収束点まで実行してから，もう一方を始める．もう一方の経路のスレッドが
  書く値を待ってスピンするスレッドは，自分の経路が先にスケジュール
  されるとハングする．書き手は，スピンしている経路が再収束点に達した後に
  しか実行されないが，その経路は決してそこに達しないからである．
  これが\emph{SIMT デッドロック}である．たとえば，レーン~0 が
  \texttt{then} 経路を取り，メモリ上のフラグがセットされるまでループし，
  レーン~1 は \texttt{else} 経路を取ってそのフラグをセットするとする．
  スタックが \texttt{then} 経路を先に実行すると，レーン~0 は永遠に
  ループし，\texttt{else} 経路は決してフェッチされない．1つのワープの
  複数のレーンが取るスピンロックも同じように失敗する．A18 ch.~3 は
  この問題を，スレッドごとにプログラムカウンタを持ち，スケジューラが
  もう一方の経路に切り替えられるようにするスタックレスの設計と合わせて
  論じている．\margin{Volta 以降の NVIDIA GPU はスレッドごとに
  プログラムカウンタと呼び出しスタックを持ち（独立スレッドスケジューリング，
  第5章），正しさについて古典的なスタックにもはや依存しない（CPG §4.1
  SIMT Architecture \cite{nvidia2024cudaguide}）．}本書の以降の部分では，
  本節のワープごとのスタックのモデルを用いる．
\end{note}

\subsection{ラージワープ}

ワープサイズはマイクロアーキテクチャ上の選択であって，法則ではない．
1回のフェッチとデコードで32，64，128のレーンを駆動でき，非常に幅の広い
ワープは SM を，レーンあたりのフロントエンドのコストが小さい幅広い
ベクトル機械に変える．そのコストはダイバージェンスである．ワープが
保持するスレッドが多いほど，そのうちの1つが他と一致しない可能性は高く，
ダイバージェンスのたびに経路が逐次化され SIMT スタックが占有される．
適切な幅は，アプリケーションの制御フローがどれだけ一様かによる．\caution{実際の
分岐結果には相関がある（\texttt{threadIdx.x / 32} に対する条件は構成上
ワープごとに一様である）ので，独立な結果は悲観的なモデルである．
これは傾向を示すものであって，実際のダイバージェンス率ではない．}

\begin{example}
  \label{ex:l06-large-warp}
  各スレッドが独立に確率 $p = 0.1$ で分岐を取るとする．$W$ スレッドの
  ワープが一様になるのは全スレッドが一致するときだけであり，その確率は
  $0.9^W + 0.1^W$ である．これは $W = 4$ で約 $0.66$，$W = 32$ で $0.034$，
  $W = 64$ で $0.0012$ となる．このモデルでは，32スレッドのワープの
  ほとんどすべての分岐がダイバージェントになる．
\end{example}

\begin{definition}[ラージワープ]
  \label{def:l06-large-warp}
  \term[らーじわーぷ]{ラージワープ}[large warp]とは，従来の32より多い
  （64，128またはそれ以上の）スレッドを持つワープであり，その命令は
  全スレッドに対して1回だけフェッチ，デコードされるので，フロントエンドの
  コストがより多くのレーンで分担される．
\end{definition}

\caution{講義 vs 論文：本章は講義に従い，ラージワープを幅の広いワープとして
扱う．Narasiman ら \cite{narasiman2011twolevel} はラージワープの
\emph{アクティブな}スレッドを毎サイクル SIMD 幅のサブワープに詰め直すので，
ダイバージェンスによって遊ぶレーンが少なくなる．2レベルのワープ
スケジューリング（第7章）と組み合わせて，従来の GPU コアに対して
19.1\,\% の向上を報告している．}%
したがって論文のラージワープは，ダイバージェンスに対抗する方法でもあり，
次のアイデアにつながる．

\subsection{動的ワープ形成}

ダイバージェント分岐の後，各経路はレーンの一部だけを働かせて実行される．
\cref{fig:l06-divergent-cfg}では BB2 と BB3 でレーンの半分が遊ぶ．
働いているレーンの割合がこの損失の尺度である．

\begin{definition}[SIMD 効率]
  \label{def:l06-simd-efficiency}
  実行の\term[しむどこうりつ]{SIMD効率}[SIMD efficiency]とは，仕事をする
  レーンスロットの数を，発行されたレーンスロットの数（ワープサイズ
  $\times$ 発行されたワープ命令数）で割ったものである．ダイバージェンスの
  ないプログラムの SIMD 効率は 100\,\% であり，
  \cref{ex:l06-if-conversion}のダイバージェントなワープでは
  $40/64 = 62.5\,\%$ である．
\end{definition}
\margin{PMPP ch.~10 \cite{hwu2022pmpp} は同じ尺度をソフトウェアに
適用する．並列リダクションを書き換えて，アクティブなスレッドを少数の
ワープに詰めたままにし，SIMD 効率を高く保つ．}

SM の多くのワープはしばしばプログラムの同じ地点にあり，同じダイバージェント
分岐の後，異なるワープが相補的なマスクを持って同じ PC で待っている．
それらのスレッドは結合できる．\margin{Fung らは，現実的な DWF の実装で，
推定 4.7\,\% の面積増加に対して平均 20.7\,\% の高速化を報告している
\cite{fung2007dwf}．}

\begin{definition}[動的ワープ形成]
  \label{def:l06-dwf}
  \term[どうてきわーぷけいせい]{動的ワープ形成}[dynamic warp formation (DWF)]\index{DWF|see{動的ワープ形成}}（DWF）
  と呼ばれる手法は，同じ PC を実行しようとしている異なるワープの
  スレッドを，実行時に，より満たされたアクティブマスクを持つ新しい
  ワープにまとめ直す（\cref{fig:l06-dwf}）．
\end{definition}

\begin{figure}[tb]
  \centering
  % Dynamic warp formation (Lesson 6).  Two warps of four threads (T0-T3,
% T4-T7) diverge on the same branch: at BB2 warp 0 has lanes 0-1 active
% and warp 1 lanes 2-3, so the two half warps fill one new warp without
% moving any thread out of its home lane; BB3 likewise.
% Redraws the outline figure m6l1_dynamicwarp.
% Usage: % Dynamic warp formation (Lesson 6).  Two warps of four threads (T0-T3,
% T4-T7) diverge on the same branch: at BB2 warp 0 has lanes 0-1 active
% and warp 1 lanes 2-3, so the two half warps fill one new warp without
% moving any thread out of its home lane; BB3 likewise.
% Redraws the outline figure m6l1_dynamicwarp.
% Usage: % Dynamic warp formation (Lesson 6).  Two warps of four threads (T0-T3,
% T4-T7) diverge on the same branch: at BB2 warp 0 has lanes 0-1 active
% and warp 1 lanes 2-3, so the two half warps fill one new warp without
% moving any thread out of its home lane; BB3 likewise.
% Redraws the outline figure m6l1_dynamicwarp.
% Usage: \input{figures/l06-dwf}
% Width: body (120 mm); the drawing is about 114 mm wide.
\begin{tikzpicture}[x=1mm, y=1mm, font=\scriptsize\sffamily,
    >={Stealth[length=1.6mm]}]
  % \lanes{x}{y}{labels}{fill}: four lane cells, "-" = idle
  \def\dashmark{-}
  \newcommand{\lanes}[4]{%
    \foreach \t [count=\l from 0] in {#3} {
      \ifx\t\dashmark
        \draw[ln-rule, dashed] (#1+\l*7,#2) rectangle ++(6,5);
      \else
        \draw[fill=#4] (#1+\l*7,#2) rectangle ++(6,5) node[midway] {\t};
      \fi
    }}
  % column heads
  \node[font=\scriptsize\sffamily\bfseries] at (13.5,30) {\iflangja{ワープ 0}{warp 0}};
  \node[font=\scriptsize\sffamily\bfseries] at (49.5,30) {\iflangja{ワープ 1}{warp 1}};
  \node[font=\scriptsize\sffamily\bfseries, align=center] at (95.5,30)
    {\iflangja{動的に形成したワープ}{dynamically formed warp}};
  \foreach \l in {0,...,3} {
    \node[text=ln-muted] at (3+\l*7,25.5) {\l};
    \node[text=ln-muted] at (39+\l*7,25.5) {\l};
    \node[text=ln-muted] at (85+\l*7,25.5) {\l};
  }
  \node[anchor=east, text=ln-muted] at (-1,25.5) {\iflangja{レーン}{lane}};
  % BB2 row
  \node[anchor=east] at (-1,19.5) {BB2};
  \lanes{0}{17}{T0,T1,-,-}{ln-source!25}
  \lanes{36}{17}{-,-,T6,T7}{ln-source!25}
  \lanes{82}{17}{T0,T1,T6,T7}{ln-source!25}
  % BB3 row
  \node[anchor=east] at (-1,9.5) {BB3};
  \lanes{0}{7}{-,-,T2,T3}{ln-tip!25}
  \lanes{36}{7}{T4,T5,-,-}{ln-tip!25}
  \lanes{82}{7}{T4,T5,T2,T3}{ln-tip!25}
  \draw[->, very thick, ln-accent] (66,14.5) -- (79,14.5);
  \node[anchor=north, text=ln-muted, align=center] at (72.5,12.5)
    {\iflangja{同じ PC で\\再グループ化}{regroup\\at equal PC}};
  \node[anchor=north west, text=ln-muted] at (0,4)
    {\iflangja{SIMD 効率 50\,\% $\times$ 4 命令スロット}{SIMD efficiency 50\,\% in 4 issue slots}};
  \node[anchor=north west, text=ln-muted] at (82,4)
    {\iflangja{100\,\% $\times$ 2 スロット}{100\,\% in 2 issue slots}};
\end{tikzpicture}

% Width: body (120 mm); the drawing is about 114 mm wide.
\begin{tikzpicture}[x=1mm, y=1mm, font=\scriptsize\sffamily,
    >={Stealth[length=1.6mm]}]
  % \lanes{x}{y}{labels}{fill}: four lane cells, "-" = idle
  \def\dashmark{-}
  \newcommand{\lanes}[4]{%
    \foreach \t [count=\l from 0] in {#3} {
      \ifx\t\dashmark
        \draw[ln-rule, dashed] (#1+\l*7,#2) rectangle ++(6,5);
      \else
        \draw[fill=#4] (#1+\l*7,#2) rectangle ++(6,5) node[midway] {\t};
      \fi
    }}
  % column heads
  \node[font=\scriptsize\sffamily\bfseries] at (13.5,30) {\iflangja{ワープ 0}{warp 0}};
  \node[font=\scriptsize\sffamily\bfseries] at (49.5,30) {\iflangja{ワープ 1}{warp 1}};
  \node[font=\scriptsize\sffamily\bfseries, align=center] at (95.5,30)
    {\iflangja{動的に形成したワープ}{dynamically formed warp}};
  \foreach \l in {0,...,3} {
    \node[text=ln-muted] at (3+\l*7,25.5) {\l};
    \node[text=ln-muted] at (39+\l*7,25.5) {\l};
    \node[text=ln-muted] at (85+\l*7,25.5) {\l};
  }
  \node[anchor=east, text=ln-muted] at (-1,25.5) {\iflangja{レーン}{lane}};
  % BB2 row
  \node[anchor=east] at (-1,19.5) {BB2};
  \lanes{0}{17}{T0,T1,-,-}{ln-source!25}
  \lanes{36}{17}{-,-,T6,T7}{ln-source!25}
  \lanes{82}{17}{T0,T1,T6,T7}{ln-source!25}
  % BB3 row
  \node[anchor=east] at (-1,9.5) {BB3};
  \lanes{0}{7}{-,-,T2,T3}{ln-tip!25}
  \lanes{36}{7}{T4,T5,-,-}{ln-tip!25}
  \lanes{82}{7}{T4,T5,T2,T3}{ln-tip!25}
  \draw[->, very thick, ln-accent] (66,14.5) -- (79,14.5);
  \node[anchor=north, text=ln-muted, align=center] at (72.5,12.5)
    {\iflangja{同じ PC で\\再グループ化}{regroup\\at equal PC}};
  \node[anchor=north west, text=ln-muted] at (0,4)
    {\iflangja{SIMD 効率 50\,\% $\times$ 4 命令スロット}{SIMD efficiency 50\,\% in 4 issue slots}};
  \node[anchor=north west, text=ln-muted] at (82,4)
    {\iflangja{100\,\% $\times$ 2 スロット}{100\,\% in 2 issue slots}};
\end{tikzpicture}

% Width: body (120 mm); the drawing is about 114 mm wide.
\begin{tikzpicture}[x=1mm, y=1mm, font=\scriptsize\sffamily,
    >={Stealth[length=1.6mm]}]
  % \lanes{x}{y}{labels}{fill}: four lane cells, "-" = idle
  \def\dashmark{-}
  \newcommand{\lanes}[4]{%
    \foreach \t [count=\l from 0] in {#3} {
      \ifx\t\dashmark
        \draw[ln-rule, dashed] (#1+\l*7,#2) rectangle ++(6,5);
      \else
        \draw[fill=#4] (#1+\l*7,#2) rectangle ++(6,5) node[midway] {\t};
      \fi
    }}
  % column heads
  \node[font=\scriptsize\sffamily\bfseries] at (13.5,30) {\iflangja{ワープ 0}{warp 0}};
  \node[font=\scriptsize\sffamily\bfseries] at (49.5,30) {\iflangja{ワープ 1}{warp 1}};
  \node[font=\scriptsize\sffamily\bfseries, align=center] at (95.5,30)
    {\iflangja{動的に形成したワープ}{dynamically formed warp}};
  \foreach \l in {0,...,3} {
    \node[text=ln-muted] at (3+\l*7,25.5) {\l};
    \node[text=ln-muted] at (39+\l*7,25.5) {\l};
    \node[text=ln-muted] at (85+\l*7,25.5) {\l};
  }
  \node[anchor=east, text=ln-muted] at (-1,25.5) {\iflangja{レーン}{lane}};
  % BB2 row
  \node[anchor=east] at (-1,19.5) {BB2};
  \lanes{0}{17}{T0,T1,-,-}{ln-source!25}
  \lanes{36}{17}{-,-,T6,T7}{ln-source!25}
  \lanes{82}{17}{T0,T1,T6,T7}{ln-source!25}
  % BB3 row
  \node[anchor=east] at (-1,9.5) {BB3};
  \lanes{0}{7}{-,-,T2,T3}{ln-tip!25}
  \lanes{36}{7}{T4,T5,-,-}{ln-tip!25}
  \lanes{82}{7}{T4,T5,T2,T3}{ln-tip!25}
  \draw[->, very thick, ln-accent] (66,14.5) -- (79,14.5);
  \node[anchor=north, text=ln-muted, align=center] at (72.5,12.5)
    {\iflangja{同じ PC で\\再グループ化}{regroup\\at equal PC}};
  \node[anchor=north west, text=ln-muted] at (0,4)
    {\iflangja{SIMD 効率 50\,\% $\times$ 4 命令スロット}{SIMD efficiency 50\,\% in 4 issue slots}};
  \node[anchor=north west, text=ln-muted] at (82,4)
    {\iflangja{100\,\% $\times$ 2 スロット}{100\,\% in 2 issue slots}};
\end{tikzpicture}

  \caption{動的ワープ形成．ワープ0と1が同じ分岐でダイバージェントになる．
    BB2 ではワープ~0 のアクティブなスレッドはレーン 0--1 に，ワープ~1 の
    ものはレーン 2--3 にあるので，レーンを離れることなく1つの完全な
    ワープを形成する．BB3 も同様にマージされる．}
  \label{fig:l06-dwf}
\end{figure}

制約はレジスタファイルである．第4章ではスレッド~$k$ のレジスタをバンク~$k$
に割り当て，レーン~$k$ が自分のバンクから読むようにした．別のレーンに
移されたスレッドは，別のレーンに属するバンクからレジスタを読む必要が
あり，このレイアウトが避けるように設計されたレジスタバンク競合を
生む．そのため DWF は各スレッドをホームレーンに置いたまま，アクティブな
レーンが重ならないワープだけをマージする（\cref{fig:l06-dwf}）．\tip{DWF が
同じ PC の2つのワープをマージできるのは，それらのマスクのビットごとの
AND が0のときだけである．新しいマスクはそれらの OR である．}
どちらもレーン 0--1 を使う2つの半ワープはマージできない．

\begin{example}
  \label{ex:l06-dwf}
  \cref{ex:l06-if-conversion}のブロックサイズを持つ分岐を実行する
  2つのワープを考え，条件がデータに依存して，\cref{fig:l06-dwf}のように
  ワープ~0 はレーン 0--1 で，ワープ~1 はレーン 2--3 で BB2 を取るとする．
  DWF がなければ BB2 と BB3 は SIMD 効率 50\,\% で
  $2 \times (6+6) = 24$ 命令スロットかかる．DWF があれば，マージされた
  ワープは 100\,\% で $6+6 = 12$ 命令スロットで済む．カーネル全体では
  $2 \times 16 = 32$ から $2\cdot2 + 12 + 2\cdot2 = 20$ ワープ命令に減り，
  SIMD 効率は 62.5\,\% から $80/80 = 100\,\%$ に上がる．
\end{example}

\begin{keypoint}
  ダイバージェント分岐は，ワープに両方の経路を部分的なマスクで実行させる．
  if変換はそれを常に両方の経路をフェッチすることと引き換えにし，SIMT
  スタックは経路を順に実行して直接後支配節点で再収束させる．幅の広い
  ワープはフロントエンドを分担するがダイバージェンスが増え（Narasiman ら
  のラージワープのマイクロアーキテクチャはアクティブなスレッドを
  サブワープに詰め直す），動的ワープ形成は各スレッドをホームレーンに
  置いたまま遊んでいるレーンを取り戻す．
\end{keypoint}

\section{レジスタの最適化}
\label{sec:l06-registers}

SM のレジスタファイルは常駐するすべてのスレッドに専用のレジスタの組を
保持する（第4章）ので，その大きさはスレッド数とスレッドあたりの
レジスタ数の積であり，世代を追って SM あたり \SI{64}{\kibi\byte}，
\SI{128}{\kibi\byte}，\SI{256}{\kibi\byte} となってきた．\source{M6，
レジスタの最適化；K12 §1.4.1 GPU Pipeline；
\cite{gebhart2011energy,gebhart2011compile,jeon2015rfv}．}
大きな SRAM は遅く電力を消費するが，GPU は遅いものを受け入れられない．
レジスタファイルはレーンが消費する速さでオペランドを供給しなければ
ならず，その読み出し帯域は実行帯域に等しくなければならない．\margin{SM
あたり \SI{256}{\kibi\byte}（V100，A100，H100）．GPU あたりでは A100 で
\SI{27648}{\kibi\byte}，H100 SXM5 で \SI{33792}{\kibi\byte}（暫定値）
\cite[Table~3]{nvidia2022hopper}．}CPU の設計から2つの対策が得られる．
よくあるアクセスを速くするための階層と，必要な記憶容量を減らすための
共有である．GPU プログラムがレジスタをどう使うかについての4つの観察が，
それらがどこに当てはまるかを示す．

\subsection{4つの観察}

\begin{description}
  \item[すべてのスレッドがアクティブなわけではない．]  同じカーネルの
    スレッドは同じレジスタ名を使うが，同時に終わるわけではない．実行は
    非同期であり，終了したスレッドも，そのワープまたはブロックが解放される
    までレジスタを占有し続ける．ダイバージェント分岐では，2つの経路が
    異なるレジスタを使うこともある．
  \item[すべてのレジスタが生存しているわけではない．]
    \texttt{add r1, r2, \#1} の後に \texttt{mov r2, \#10} が続くとき，
    \texttt{r2} の値は \texttt{add} の後二度と読まれない．\texttt{mov} が
    上書きするまで，それは非生存である．
  \item[生産者と消費者は近い．]
    \texttt{mov r2, \#10} の後に \texttt{add r3, r2, \#20} が続くとき，
    値はすぐ次の命令によって消費され，スケジューラは2つを近くに保てる．
    そのような値は大きなレジスタファイルに届く必要すらないかもしれない．
  \item[ほとんどの値は一度しか読まれない．]
    \texttt{mov r2, \#10} の後に \texttt{add r2, r2, \#20} が続くとき，
    \texttt{r2} の最初の値は1つの命令に読まれた後に上書きされる．
    多くの値がこのパターンに従い，それらを主レジスタファイルに格納する
    ことには疑問がある．\caution{講義 vs 論文：講義はレジスタの値の
    ほぼ 70\,\% が一度しか読まれないとする．Gebhart らは自分たちの
    ワークロード全体で\enquote{up to 70\,\%}と測定しており，HPC のコードは
    より多く再利用する（値の 40\,\% が2回以上読まれる）．また，一度だけ
    \emph{かつ}生成から3命令以内に読まれるのは\enquote{up to 40\,\%}
    である \cite{gebhart2011energy}．}
\end{description}

\begin{definition}[生存レジスタと非生存レジスタ]
  \label{def:l06-live-dead}
  プログラムのある地点でレジスタが
  \term[せいぞんれじすた]{生存レジスタ}[live register]であるとは，
  それが保持する値が上書きされる前に後で読まれることであり，そうでなければ
  \term[ひせいぞんれじすた]{非生存レジスタ}[dead register]である．
  値の\term[あたいのせいぞんきかん]{値の生存期間}[value lifetime]とは，
  それを生成する命令から最後に読まれるまでの区間である．
\end{definition}

2番目の観察は CPU より GPU で重い（\cref{fig:l06-live-dead}）．
CPU では \texttt{add} と \texttt{mov} は連続するサイクルで実行され，
\texttt{r2} が非生存なのは1サイクル以下である．GPU ではスケジューラが
ワープをインタリーブするので，同じワープの次の命令は数十ワープ命令後に
発行されることがある．その間ずっと，\texttt{r2} の非生存の値がワープの
全スレッドのレジスタを1つずつ占有する．\margin{このインタリーブを
決めるワープスケジューラは第7章で論じる．貪欲なスケジューラは間隔を
縮め，ラウンドロビンのスケジューラは間隔を準備のできたワープの数まで
引き延ばす．}

\begin{figure}[tb]
  \centering
  % Live and dead registers on a CPU and on a GPU (Lesson 6).  After
% "add r1, r2, #1" register r2 is dead until "mov r2, #10" overwrites it.
% On a CPU the two instructions issue in consecutive cycles; on a GPU the
% issue slots in between go to other warps, so r2 of warp w1 holds a
% dead value for tens of cycles.
% Redraws the outline figure m6l2_o2.
% Usage: % Live and dead registers on a CPU and on a GPU (Lesson 6).  After
% "add r1, r2, #1" register r2 is dead until "mov r2, #10" overwrites it.
% On a CPU the two instructions issue in consecutive cycles; on a GPU the
% issue slots in between go to other warps, so r2 of warp w1 holds a
% dead value for tens of cycles.
% Redraws the outline figure m6l2_o2.
% Usage: % Live and dead registers on a CPU and on a GPU (Lesson 6).  After
% "add r1, r2, #1" register r2 is dead until "mov r2, #10" overwrites it.
% On a CPU the two instructions issue in consecutive cycles; on a GPU the
% issue slots in between go to other warps, so r2 of warp w1 holds a
% dead value for tens of cycles.
% Redraws the outline figure m6l2_o2.
% Usage: \input{figures/l06-live-dead}
% Width: body (120 mm); the drawing is about 116 mm wide.
\begin{tikzpicture}[x=1mm, y=1mm, font=\scriptsize\sffamily,
    slot/.style={draw, minimum width=14mm, minimum height=6mm,
                 inner sep=1pt, align=center}]
  % CPU
  \node[anchor=west, font=\scriptsize\sffamily\bfseries] at (0,31) {(a) CPU};
  \node[slot, fill=ln-accent!20] at (24,31) {\texttt{add r1,r2,\#1}};
  \node[slot, fill=ln-tip!25] at (40,31) {\texttt{mov r2,\#10}};
  \draw[ln-caution, thick] (31,26.5) -- ++(0,-1) -- ++(2,0) -- ++(0,1);
  \node[anchor=west, text=ln-caution] at (34,24.5)
    {\iflangja{\texttt{r2} 非生存：1 サイクル未満}{\texttt{r2} dead: under one cycle}};
  % GPU
  \node[anchor=west, font=\scriptsize\sffamily\bfseries] at (0,13) {(b) GPU};
  \foreach \t/\w/\c [count=\i from 0] in
      {add/w1/ln-accent!20, add/w2/ln-box, add/w3/ln-box, add/w4/ln-box,
       {\ldots}/{}/white, mov/w1/ln-tip!25} {
    \node[slot, fill=\c, minimum width=13mm] at (24+\i*15,13) {\texttt{\t}\\\w};
  }
  \draw[ln-caution, thick] (31,8.5) -- ++(0,-2) -- ++(72,0) -- ++(0,2);
  \node[anchor=north, text=ln-caution] at (67,6)
    {\iflangja{w1 の \texttt{r2} は非生存（数十命令スロットの間）}{\texttt{r2} of w1 dead for tens of issue slots}};
  \node[anchor=north, text=ln-muted] at (67,20.5)
    {\iflangja{1 命令スロットごとに 1 ワープ命令}{one warp instruction per issue slot}};
\end{tikzpicture}

% Width: body (120 mm); the drawing is about 116 mm wide.
\begin{tikzpicture}[x=1mm, y=1mm, font=\scriptsize\sffamily,
    slot/.style={draw, minimum width=14mm, minimum height=6mm,
                 inner sep=1pt, align=center}]
  % CPU
  \node[anchor=west, font=\scriptsize\sffamily\bfseries] at (0,31) {(a) CPU};
  \node[slot, fill=ln-accent!20] at (24,31) {\texttt{add r1,r2,\#1}};
  \node[slot, fill=ln-tip!25] at (40,31) {\texttt{mov r2,\#10}};
  \draw[ln-caution, thick] (31,26.5) -- ++(0,-1) -- ++(2,0) -- ++(0,1);
  \node[anchor=west, text=ln-caution] at (34,24.5)
    {\iflangja{\texttt{r2} 非生存：1 サイクル未満}{\texttt{r2} dead: under one cycle}};
  % GPU
  \node[anchor=west, font=\scriptsize\sffamily\bfseries] at (0,13) {(b) GPU};
  \foreach \t/\w/\c [count=\i from 0] in
      {add/w1/ln-accent!20, add/w2/ln-box, add/w3/ln-box, add/w4/ln-box,
       {\ldots}/{}/white, mov/w1/ln-tip!25} {
    \node[slot, fill=\c, minimum width=13mm] at (24+\i*15,13) {\texttt{\t}\\\w};
  }
  \draw[ln-caution, thick] (31,8.5) -- ++(0,-2) -- ++(72,0) -- ++(0,2);
  \node[anchor=north, text=ln-caution] at (67,6)
    {\iflangja{w1 の \texttt{r2} は非生存（数十命令スロットの間）}{\texttt{r2} of w1 dead for tens of issue slots}};
  \node[anchor=north, text=ln-muted] at (67,20.5)
    {\iflangja{1 命令スロットごとに 1 ワープ命令}{one warp instruction per issue slot}};
\end{tikzpicture}

% Width: body (120 mm); the drawing is about 116 mm wide.
\begin{tikzpicture}[x=1mm, y=1mm, font=\scriptsize\sffamily,
    slot/.style={draw, minimum width=14mm, minimum height=6mm,
                 inner sep=1pt, align=center}]
  % CPU
  \node[anchor=west, font=\scriptsize\sffamily\bfseries] at (0,31) {(a) CPU};
  \node[slot, fill=ln-accent!20] at (24,31) {\texttt{add r1,r2,\#1}};
  \node[slot, fill=ln-tip!25] at (40,31) {\texttt{mov r2,\#10}};
  \draw[ln-caution, thick] (31,26.5) -- ++(0,-1) -- ++(2,0) -- ++(0,1);
  \node[anchor=west, text=ln-caution] at (34,24.5)
    {\iflangja{\texttt{r2} 非生存：1 サイクル未満}{\texttt{r2} dead: under one cycle}};
  % GPU
  \node[anchor=west, font=\scriptsize\sffamily\bfseries] at (0,13) {(b) GPU};
  \foreach \t/\w/\c [count=\i from 0] in
      {add/w1/ln-accent!20, add/w2/ln-box, add/w3/ln-box, add/w4/ln-box,
       {\ldots}/{}/white, mov/w1/ln-tip!25} {
    \node[slot, fill=\c, minimum width=13mm] at (24+\i*15,13) {\texttt{\t}\\\w};
  }
  \draw[ln-caution, thick] (31,8.5) -- ++(0,-2) -- ++(72,0) -- ++(0,2);
  \node[anchor=north, text=ln-caution] at (67,6)
    {\iflangja{w1 の \texttt{r2} は非生存（数十命令スロットの間）}{\texttt{r2} of w1 dead for tens of issue slots}};
  \node[anchor=north, text=ln-muted] at (67,20.5)
    {\iflangja{1 命令スロットごとに 1 ワープ命令}{one warp instruction per issue slot}};
\end{tikzpicture}

  \caption{\texttt{add r1,r2,\#1} の後，レジスタ \texttt{r2} は
    \texttt{mov r2,\#10} まで非生存である．(a) CPU では2つの命令は
    隣接する．(b) GPU では間の命令スロットは他のワープのものなので，
    ワープ \texttt{w1} の \texttt{r2} の非生存区間は数十命令にわたる．}
  \label{fig:l06-live-dead}
\end{figure}


\begin{example}
  \label{ex:l06-live-dead}
  SM が24ワープをラウンドロビンでインタリーブし，1巡ごとに各ワープ
  1命令を発行するので，1つのワープの連続する命令は24サイクル離れる．
  命令 $k$ が生成し，命令 $k+1$ が最後に読み，命令 $k+5$ が上書きする値は，
  $24$ サイクル生存し $4 \times 24 = 96$ サイクル非生存である．その
  レジスタが有用なデータを保持するのは時間の $24/120 = 20\,\%$ である．
  1サイクルに1命令を発行する CPU では，同じ値は1サイクル生存し
  4サイクル非生存であり，同じく 20\,\% であるが，非生存レジスタのコストは
  小さなファイルの1エントリである．GPU では，数百キビバイトのファイルで
  スレッドごとに1エントリを96サイクル占める．
\end{example}

\subsection{階層型レジスタファイル}

3番目と4番目の観察は階層を示唆する．生成後すぐに，しかも一度だけ読まれる
値は，演算器の近くにある小さな構造から供給できる．\tip{生産者の直後に
発行される命令によって一度だけ読まれる値は，理想的なキャッシュ候補である．
それは小さなファイルで読まれて死に，主レジスタファイルへの書き込みを
一度も生じない．}

\begin{definition}[階層型レジスタファイル]
  \label{def:l06-hierarchical-rf}
  \term[かいそうがたれじすたふぁいる]{階層型レジスタファイル}[hierarchical register file]
  は，より小さく速いレジスタ記憶を
  \term[しゅれじすたふぁいる]{主レジスタファイル}[main register file (MRF)]\index{MRF|see{主レジスタファイル}}（MRF）
  の前に置く．\cref{fig:l06-rf-hierarchy}の2つの形態は，最近生成された
  値を保持して MRF への読み書きをふるい分ける1段または2段の小さな階層
  （L0 と L1 RF）である
  \term[れじすたふぁいるきゃっしゅ]{レジスタファイルキャッシュ}[register file cache (RFC)]\index{RFC|see{レジスタファイルキャッシュ}}（RFC）
  と，高速部と低速部を持ち，低速で大きい部分をより高密度の技術で作れる
  \emph{分割型}レジスタファイルである．
\end{definition}

\begin{figure}[tb]
  \centering
  % Two hierarchical register files (Lesson 6).  (a) A register file
% cache: small L0/L1 register files in front of the main register file,
% closest to the ALUs.  (b) A partitioned register file: a fast SRAM
% segment and a slow, large segment that may be built in non-volatile
% memory.
% Redraws the outline figure m6l2_opt1.
% Usage: % Two hierarchical register files (Lesson 6).  (a) A register file
% cache: small L0/L1 register files in front of the main register file,
% closest to the ALUs.  (b) A partitioned register file: a fast SRAM
% segment and a slow, large segment that may be built in non-volatile
% memory.
% Redraws the outline figure m6l2_opt1.
% Usage: % Two hierarchical register files (Lesson 6).  (a) A register file
% cache: small L0/L1 register files in front of the main register file,
% closest to the ALUs.  (b) A partitioned register file: a fast SRAM
% segment and a slow, large segment that may be built in non-volatile
% memory.
% Redraws the outline figure m6l2_opt1.
% Usage: \input{figures/l06-rf-hierarchy}
% Width: body (120 mm); the drawing is about 110 mm wide.
\begin{tikzpicture}[x=1mm, y=1mm, font=\scriptsize\sffamily,
    >={Stealth[length=1.6mm]}]
  \node[anchor=west, font=\scriptsize\sffamily\bfseries] at (0,40)
    {\iflangja{(a) レジスタファイルキャッシュ}{(a) register file cache}};
  \draw[fill=ln-box] (14,30) rectangle ++(20,6) node[midway] {\iflangja{演算器}{ALUs}};
  \draw[fill=ln-accent!35] (16,22) rectangle ++(16,6) node[midway] {L0 RF};
  \draw[fill=ln-accent!22] (10,14) rectangle ++(28,6) node[midway] {L1 RF};
  \draw[fill=ln-accent!10] (2,4) rectangle ++(44,8)
    node[midway] {\iflangja{主レジスタファイル (MRF)}{main register file (MRF)}};
  \draw[<->] (24,28) -- (24,30);
  \draw[<->] (24,20) -- (24,22);
  \draw[<->] (24,12) -- (24,14);
  \draw[->, ln-muted] (50,34) -- (50,6);
  \node[anchor=west, text=ln-muted] at (51,33) {\iflangja{小・高速}{small, fast}};
  \node[anchor=west, text=ln-muted] at (51,7) {\iflangja{大・低速}{large, slow}};
  % partitioned
  \node[anchor=west, font=\scriptsize\sffamily\bfseries] at (68,40)
    {\iflangja{(b) 分割型レジスタファイル}{(b) partitioned register file}};
  \draw[fill=ln-box] (76,30) rectangle ++(28,6) node[midway] {\iflangja{演算器}{ALUs}};
  \draw[fill=ln-accent!30] (70,14) rectangle ++(16,12)
    node[midway, align=center] {\iflangja{高速部\\(SRAM)}{fast\\segment\\(SRAM)}};
  \draw[fill=ln-tip!20] (86,14) rectangle ++(26,12)
    node[midway, align=center] {\iflangja{低速・大容量部\\(NVM も可)}{slow, large\\segment\\(may be NVM)}};
  \draw[<->] (82,26) -- (82,30);
  \draw[<->] (98,26) -- (98,30);
  \node[anchor=north, text=ln-muted, align=center] at (91,11)
    {\iflangja{よく使う値を高速部に置く}{frequently used values\\in the fast segment}};
\end{tikzpicture}

% Width: body (120 mm); the drawing is about 110 mm wide.
\begin{tikzpicture}[x=1mm, y=1mm, font=\scriptsize\sffamily,
    >={Stealth[length=1.6mm]}]
  \node[anchor=west, font=\scriptsize\sffamily\bfseries] at (0,40)
    {\iflangja{(a) レジスタファイルキャッシュ}{(a) register file cache}};
  \draw[fill=ln-box] (14,30) rectangle ++(20,6) node[midway] {\iflangja{演算器}{ALUs}};
  \draw[fill=ln-accent!35] (16,22) rectangle ++(16,6) node[midway] {L0 RF};
  \draw[fill=ln-accent!22] (10,14) rectangle ++(28,6) node[midway] {L1 RF};
  \draw[fill=ln-accent!10] (2,4) rectangle ++(44,8)
    node[midway] {\iflangja{主レジスタファイル (MRF)}{main register file (MRF)}};
  \draw[<->] (24,28) -- (24,30);
  \draw[<->] (24,20) -- (24,22);
  \draw[<->] (24,12) -- (24,14);
  \draw[->, ln-muted] (50,34) -- (50,6);
  \node[anchor=west, text=ln-muted] at (51,33) {\iflangja{小・高速}{small, fast}};
  \node[anchor=west, text=ln-muted] at (51,7) {\iflangja{大・低速}{large, slow}};
  % partitioned
  \node[anchor=west, font=\scriptsize\sffamily\bfseries] at (68,40)
    {\iflangja{(b) 分割型レジスタファイル}{(b) partitioned register file}};
  \draw[fill=ln-box] (76,30) rectangle ++(28,6) node[midway] {\iflangja{演算器}{ALUs}};
  \draw[fill=ln-accent!30] (70,14) rectangle ++(16,12)
    node[midway, align=center] {\iflangja{高速部\\(SRAM)}{fast\\segment\\(SRAM)}};
  \draw[fill=ln-tip!20] (86,14) rectangle ++(26,12)
    node[midway, align=center] {\iflangja{低速・大容量部\\(NVM も可)}{slow, large\\segment\\(may be NVM)}};
  \draw[<->] (82,26) -- (82,30);
  \draw[<->] (98,26) -- (98,30);
  \node[anchor=north, text=ln-muted, align=center] at (91,11)
    {\iflangja{よく使う値を高速部に置く}{frequently used values\\in the fast segment}};
\end{tikzpicture}

% Width: body (120 mm); the drawing is about 110 mm wide.
\begin{tikzpicture}[x=1mm, y=1mm, font=\scriptsize\sffamily,
    >={Stealth[length=1.6mm]}]
  \node[anchor=west, font=\scriptsize\sffamily\bfseries] at (0,40)
    {\iflangja{(a) レジスタファイルキャッシュ}{(a) register file cache}};
  \draw[fill=ln-box] (14,30) rectangle ++(20,6) node[midway] {\iflangja{演算器}{ALUs}};
  \draw[fill=ln-accent!35] (16,22) rectangle ++(16,6) node[midway] {L0 RF};
  \draw[fill=ln-accent!22] (10,14) rectangle ++(28,6) node[midway] {L1 RF};
  \draw[fill=ln-accent!10] (2,4) rectangle ++(44,8)
    node[midway] {\iflangja{主レジスタファイル (MRF)}{main register file (MRF)}};
  \draw[<->] (24,28) -- (24,30);
  \draw[<->] (24,20) -- (24,22);
  \draw[<->] (24,12) -- (24,14);
  \draw[->, ln-muted] (50,34) -- (50,6);
  \node[anchor=west, text=ln-muted] at (51,33) {\iflangja{小・高速}{small, fast}};
  \node[anchor=west, text=ln-muted] at (51,7) {\iflangja{大・低速}{large, slow}};
  % partitioned
  \node[anchor=west, font=\scriptsize\sffamily\bfseries] at (68,40)
    {\iflangja{(b) 分割型レジスタファイル}{(b) partitioned register file}};
  \draw[fill=ln-box] (76,30) rectangle ++(28,6) node[midway] {\iflangja{演算器}{ALUs}};
  \draw[fill=ln-accent!30] (70,14) rectangle ++(16,12)
    node[midway, align=center] {\iflangja{高速部\\(SRAM)}{fast\\segment\\(SRAM)}};
  \draw[fill=ln-tip!20] (86,14) rectangle ++(26,12)
    node[midway, align=center] {\iflangja{低速・大容量部\\(NVM も可)}{slow, large\\segment\\(may be NVM)}};
  \draw[<->] (82,26) -- (82,30);
  \draw[<->] (98,26) -- (98,30);
  \node[anchor=north, text=ln-muted, align=center] at (91,11)
    {\iflangja{よく使う値を高速部に置く}{frequently used values\\in the fast segment}};
\end{tikzpicture}

  \caption{2つの階層型レジスタファイル．(a) レジスタファイルキャッシュ：
    演算器と主レジスタファイルの間にある小さな L0 と L1 のレジスタ
    ファイル．(b) 分割型レジスタファイル：高速部と，不揮発性メモリを
    使いうる低速で大きい部分．}
  \label{fig:l06-rf-hierarchy}
\end{figure}

RFC があれば，ほとんどの読み出しは小さなファイルでヒットし，追い出される
前に死ぬ値は MRF に書き込む必要がない．そのため MRF を小さくするか，
同じ大きさのままアクセス頻度を下げることができる．小さな SRAM は大きな
ものよりアクセスあたりのコストが小さいので，これはエネルギーを
節約する．\margin{Gebhart ら：スレッドあたり6エントリの RFC で MRF の
読み出しは 50\,\%，書き込みは 59\,\% 減る．スレッドの4分の1をアクティブに
保つ2レベルスケジューラと組み合わせると，レジスタファイルのエネルギーは
36\,\% 下がる \cite{gebhart2011energy}．}GPU ではこれが CPU より容易である．
K12 §1.4.1 が述べるように GPU は精密例外をサポートしないので，
\enquote{register files do not have to maintain architectural states}
\cite{kim2012gpgpu}，値はキャッシュの中だけに存在してよい．

分割型ファイルの低速部の候補は
\term[ふきはつせいめもり]{不揮発性メモリ}[non-volatile memory (NVM)]\index{NVM|see{不揮発性メモリ}}（NVM），
すなわち電力なしで内容を保持するメモリである．NVM で作った低速で大きい
部分はエネルギーの節約を狙ったものであり，頻繁に使われる値は高速部に
とどまる．\margin{レジスタファイルの NVM は研究の方向であって，
出荷された設計ではない．現在の GPU はレジスタファイル全体を SRAM で
作っている．}

キャッシュの判断は，値の生存期間を知っているコンパイラに移すこともできる．
Gebhart，Keckler と Dally \cite{gebhart2011compile} はハードウェアの RFC を，
コンパイラが内容を割り当てるスレッドごとの小さなファイルである
\term[おぺらんどれじすたふぁいる]{オペランドレジスタファイル}[operand register file (ORF)]\index{ORF|see{オペランドレジスタファイル}}（ORF）
と，最新の結果を即座の再利用のために保持するスレッドあたり1エントリの
\term[らすとりざるとふぁいる]{ラストリザルトファイル}[last result file (LRF)]\index{LRF|see{ラストリザルトファイル}}（LRF）
で置き換える．\margin{コンパイラ管理の階層はレジスタファイルの
エネルギーを最大 54\,\% 削減すると報告されている．ハードウェアのみの
キャッシュでは 34\,\% である．彼らのベースラインでは，レジスタファイルは
プロセッサの動的エネルギーの 15--20\,\% を消費する
\cite{gebhart2011compile}．}コンパイラがすべての値を明示的に配置するので，
ORF にはタグもハードウェアによる置換も要らない．これはスクラッチパッドと
キャッシュを分けるのと同じ性質である（\cref{sec:l06-unified}）．

\subsection{レジスタファイル仮想化}

最初の2つの観察は共有を示唆する．どの瞬間にも多くのレジスタが非生存で
あるか，終了したスレッドに属しているなら，名前の数より少ない物理
レジスタで足りる \cite{jeon2015rfv}．

\begin{definition}[アーキテクチャレジスタと物理レジスタ]
  \label{def:l06-arch-phys}
  \term[あーきてくちゃれじすた]{アーキテクチャレジスタ}[architectural register]
  とは，プログラムから見たレジスタ名（スレッド~3 の \texttt{R1}）であり，
  \term[ぶつりれじすた]{物理レジスタ}[physical register]とは，
  レジスタファイルの記憶エントリである．実行時のアーキテクチャレジスタから
  物理レジスタへの対応付けが
  \term[れじすたりねーみんぐ]{レジスタリネーミング}[register renaming]
  である．アウトオブオーダの CPU は，偽の依存を取り除くために，少数の
  アーキテクチャレジスタを\emph{より多くの}物理レジスタにリネームする．
\end{definition}

\begin{definition}[レジスタファイル仮想化]
  \label{def:l06-rf-virtualization}
  \term[れじすたふぁいるかそうか]{レジスタファイル仮想化}[register file virtualization]
  と呼ばれる手法は，全スレッドのアーキテクチャレジスタの名前空間を保った
  まま，それを\emph{より小さな}物理レジスタファイルに対応付ける．物理
  レジスタは値が生成されるときに割り当てられ，値が死ぬときに解放されるので，
  値が異なる時刻に生存する名前どうしが記憶を共有する
  （\cref{fig:l06-rf-virtualization}）．プログラムの意味は変わらず，
  ハードウェアだけが小さくなる．
\end{definition}

\begin{figure}[tb]
  \centering
  % Register file virtualisation (Lesson 6).  Six architectural register
% names (R1 of threads 1-3, R2 of threads 4-6) whose values are live in
% disjoint intervals share two physical registers: P0 and P1.
% Redraws the outline figure m6l2_opt2 (lifetime bars added).
% Usage: % Register file virtualisation (Lesson 6).  Six architectural register
% names (R1 of threads 1-3, R2 of threads 4-6) whose values are live in
% disjoint intervals share two physical registers: P0 and P1.
% Redraws the outline figure m6l2_opt2 (lifetime bars added).
% Usage: % Register file virtualisation (Lesson 6).  Six architectural register
% names (R1 of threads 1-3, R2 of threads 4-6) whose values are live in
% disjoint intervals share two physical registers: P0 and P1.
% Redraws the outline figure m6l2_opt2 (lifetime bars added).
% Usage: \input{figures/l06-rf-virtualization}
% Width: body (120 mm); the drawing is about 112 mm wide.
\begin{tikzpicture}[x=1mm, y=1mm, font=\scriptsize\sffamily,
    >={Stealth[length=1.6mm]}]
  \node[anchor=south west, font=\scriptsize\sffamily\bfseries] at (0,40)
    {\iflangja{名前空間}{name space}};
  \node[anchor=south west, font=\scriptsize\sffamily\bfseries] at (18,40)
    {\iflangja{値が生存する区間}{interval in which the value is live}};
  \node[anchor=south west, font=\scriptsize\sffamily\bfseries] at (90,40)
    {\iflangja{物理レジスタ}{physical registers}};
  \draw[ln-rule] (18,2) -- (82,2);
  \draw[->, ln-muted] (18,0) -- (82,0) node[right] {\iflangja{時間}{time}};
  \foreach \n/\a/\b/\c [count=\i from 0] in
      {Th1.R1/0/16/ln-accent!35, Th2.R1/20/36/ln-accent!35, Th3.R1/42/60/ln-accent!35,
       Th4.R2/6/22/ln-tip!35, Th5.R2/26/42/ln-tip!35, Th6.R2/46/62/ln-tip!35} {
    \node[anchor=west, font=\scriptsize\ttfamily] at (0,36-\i*6) {\n};
    \draw[ln-rule, dotted] (18,36-\i*6) -- (82,36-\i*6);
    \draw[fill=\c] (18+\a,34.7-\i*6) rectangle (18+\b,37.3-\i*6);
  }
  \draw[fill=ln-accent!35] (92,27) rectangle ++(14,6) node[midway] {P0};
  \draw[fill=ln-tip!35] (92,9) rectangle ++(14,6) node[midway] {P1};
  \draw[ln-muted] (83,38) -- ++(1.5,0) -- ++(0,-14) -- ++(-1.5,0);
  \draw[ln-muted] (83,20) -- ++(1.5,0) -- ++(0,-14) -- ++(-1.5,0);
  \draw[->] (84.5,31) -- (92,30);
  \draw[->] (84.5,13) -- (92,12);
\end{tikzpicture}

% Width: body (120 mm); the drawing is about 112 mm wide.
\begin{tikzpicture}[x=1mm, y=1mm, font=\scriptsize\sffamily,
    >={Stealth[length=1.6mm]}]
  \node[anchor=south west, font=\scriptsize\sffamily\bfseries] at (0,40)
    {\iflangja{名前空間}{name space}};
  \node[anchor=south west, font=\scriptsize\sffamily\bfseries] at (18,40)
    {\iflangja{値が生存する区間}{interval in which the value is live}};
  \node[anchor=south west, font=\scriptsize\sffamily\bfseries] at (90,40)
    {\iflangja{物理レジスタ}{physical registers}};
  \draw[ln-rule] (18,2) -- (82,2);
  \draw[->, ln-muted] (18,0) -- (82,0) node[right] {\iflangja{時間}{time}};
  \foreach \n/\a/\b/\c [count=\i from 0] in
      {Th1.R1/0/16/ln-accent!35, Th2.R1/20/36/ln-accent!35, Th3.R1/42/60/ln-accent!35,
       Th4.R2/6/22/ln-tip!35, Th5.R2/26/42/ln-tip!35, Th6.R2/46/62/ln-tip!35} {
    \node[anchor=west, font=\scriptsize\ttfamily] at (0,36-\i*6) {\n};
    \draw[ln-rule, dotted] (18,36-\i*6) -- (82,36-\i*6);
    \draw[fill=\c] (18+\a,34.7-\i*6) rectangle (18+\b,37.3-\i*6);
  }
  \draw[fill=ln-accent!35] (92,27) rectangle ++(14,6) node[midway] {P0};
  \draw[fill=ln-tip!35] (92,9) rectangle ++(14,6) node[midway] {P1};
  \draw[ln-muted] (83,38) -- ++(1.5,0) -- ++(0,-14) -- ++(-1.5,0);
  \draw[ln-muted] (83,20) -- ++(1.5,0) -- ++(0,-14) -- ++(-1.5,0);
  \draw[->] (84.5,31) -- (92,30);
  \draw[->] (84.5,13) -- (92,12);
\end{tikzpicture}

% Width: body (120 mm); the drawing is about 112 mm wide.
\begin{tikzpicture}[x=1mm, y=1mm, font=\scriptsize\sffamily,
    >={Stealth[length=1.6mm]}]
  \node[anchor=south west, font=\scriptsize\sffamily\bfseries] at (0,40)
    {\iflangja{名前空間}{name space}};
  \node[anchor=south west, font=\scriptsize\sffamily\bfseries] at (18,40)
    {\iflangja{値が生存する区間}{interval in which the value is live}};
  \node[anchor=south west, font=\scriptsize\sffamily\bfseries] at (90,40)
    {\iflangja{物理レジスタ}{physical registers}};
  \draw[ln-rule] (18,2) -- (82,2);
  \draw[->, ln-muted] (18,0) -- (82,0) node[right] {\iflangja{時間}{time}};
  \foreach \n/\a/\b/\c [count=\i from 0] in
      {Th1.R1/0/16/ln-accent!35, Th2.R1/20/36/ln-accent!35, Th3.R1/42/60/ln-accent!35,
       Th4.R2/6/22/ln-tip!35, Th5.R2/26/42/ln-tip!35, Th6.R2/46/62/ln-tip!35} {
    \node[anchor=west, font=\scriptsize\ttfamily] at (0,36-\i*6) {\n};
    \draw[ln-rule, dotted] (18,36-\i*6) -- (82,36-\i*6);
    \draw[fill=\c] (18+\a,34.7-\i*6) rectangle (18+\b,37.3-\i*6);
  }
  \draw[fill=ln-accent!35] (92,27) rectangle ++(14,6) node[midway] {P0};
  \draw[fill=ln-tip!35] (92,9) rectangle ++(14,6) node[midway] {P1};
  \draw[ln-muted] (83,38) -- ++(1.5,0) -- ++(0,-14) -- ++(-1.5,0);
  \draw[ln-muted] (83,20) -- ++(1.5,0) -- ++(0,-14) -- ++(-1.5,0);
  \draw[->] (84.5,31) -- (92,30);
  \draw[->] (84.5,13) -- (92,12);
\end{tikzpicture}

  \caption{レジスタファイル仮想化．スレッド 1--3 の \texttt{R1} の値は
    互いに重ならない区間で生存し，物理レジスタ P0 を共有する．スレッド
    4--6 の \texttt{R2} の値は P1 を共有する．6つの名前に2つの物理
    レジスタで足りる．}
  \label{fig:l06-rf-virtualization}
\end{figure}

\caution{方向は CPU のリネーミングと逆である．CPU はアーキテクチャ
レジスタより多くの物理レジスタを持つが，GPU のレジスタファイル仮想化では
少ない．仮想化されたファイルは資源の共有であり，維持すべきなのは名前空間
だけである．}%
対応付けにはオペランドごとに表の参照がかかり，設計は値が死んだちょうど
そのときにレジスタを解放しなければならない．これにはコンパイラからの
生存期間の情報（各値の最後の読み出し）が必要である．また物理レジスタが
尽きたときの進行も保証しなければならない．すべてが使われているときに
新しいレジスタを必要とするワープは，他のワープが解放するのを待たねば
ならず，すべてのワープが待てばどれも進まない．

\begin{example}
  \label{ex:l06-rf-virtualization}
  あるカーネルがスレッドあたり32レジスタを使い，SM が1024スレッドを常駐
  させるとすると，名前空間は $32 \times 1024 = \num{32768}$ レジスタ，
  1つ \SI{4}{\byte} として \SI{128}{\kibi\byte} を持つ．プロファイリングの
  結果，非生存レジスタのためと終了したスレッドのために，どの瞬間にも
  生存する値を持つのは多くとも 60\,\% だと分かったとする．すると
  $0.6 \times \num{32768} \approx \num{19661}$ 個の物理レジスタ，約
  \SI{77}{\kibi\byte} の仮想化されたファイルですべての生存する値を保持でき，
  残りの \SI{51}{\kibi\byte} は取り除くか，常駐スレッド数を増やすために
  残しておける．\margin{計算のために選んだ数値であり，測定された生存割合を
  意味するものではない．}
\end{example}

\begin{keypoint}
  GPU のレジスタファイルは，常駐するすべてのスレッドのレジスタを保持する
  ので大きく，すべてのレーンにオペランドを供給するので速くなければならない．
  ほとんどの値は長い区間にわたって非生存であるか，生成後すぐに一度だけ
  読まれる．これが，レイテンシとエネルギーのための階層（レジスタファイル
  キャッシュ，ORF/LRF，高速部と低速部への分割）と，容量のための仮想化
  （名前より少ない物理レジスタ）を正当化する．
\end{keypoint}

\section{レジスタファイルと共有メモリの統合}
\label{sec:l06-unified}

第3章では，占有率はカーネルが最初に使い果たす資源，すなわちスレッド
スロット，レジスタ，共有メモリのいずれかによって制限されることを
示した．\source{M6，レジスタファイルと共有メモリの統合；
\cite{gebhart2012unifying}；A18 ch.~4；CPG Compute Capability 7.x--9.0，
Shared Memory．}カーネルによって後の2つの使い方は大きく異なるので，
それらの間の固定された分割は記憶を無駄にする．本節では3種類の
オンチップ記憶を比較し，統合された設計を紹介する．

\subsection{占有率再訪}

\SI{128}{\kibi\byte} のレジスタファイル（\num{32768} レジスタ）と
\SI{64}{\kibi\byte} の共有メモリを持つ SM と，\margin{構成可能な分割は
統合設計より前からある．Fermi の SM（2009年）は \SI{64}{\kibi\byte} を
\SI{48}{\kibi\byte} の共有メモリと \SI{16}{\kibi\byte} の L1，または
その逆に分割する \cite{nvidia2009fermi}．}ブロックあたり256スレッドで
起動される2つのカーネルを考える（\cref{fig:l06-unified-pool}）．
\begin{itemize}
  \item カーネル~A はスレッドあたり64レジスタを使い，共有メモリは使わない．
    ブロックあたりのレジスタは
    $64 \times 256 \times \SI{4}{\byte} = \SI{64}{\kibi\byte}$ なので，
    2ブロックでレジスタファイルが埋まり，共有メモリ全体が遊ぶ．
  \item カーネル~B はスレッドあたり16レジスタと，ブロックあたり
    \SI{16}{\kibi\byte} の共有メモリを使う．ブロックあたりのレジスタは
    \SI{16}{\kibi\byte} である．共有メモリは4ブロックを受け入れ，それらは
    レジスタファイルの半分しか使わない．\margin{常駐スレッド数の上限
（Compute Capability 7.0，8.0，9.0 では SM あたり2048，CPG Table~21
\cite{nvidia2024cudaguide}）は依然として適用される．カーネル~B の
6ブロックは1536スレッドである．}
\end{itemize}

\begin{figure}[tb]
  \centering
  % Split versus unified on-chip storage (Lesson 6, own example).  An SM
% with a 128 KiB register file and 64 KiB of shared memory, or one
% 192 KiB pool.  Kernel A: 64 KiB of registers per block, no shared
% memory.  Kernel B: 16 KiB of registers and 16 KiB of shared memory per
% block.  Each filled segment is one resident block's allocation; white
% is idle.  Scale 0.45 mm per KiB.
% Original (no outline figure for this lesson).
% Usage: % Split versus unified on-chip storage (Lesson 6, own example).  An SM
% with a 128 KiB register file and 64 KiB of shared memory, or one
% 192 KiB pool.  Kernel A: 64 KiB of registers per block, no shared
% memory.  Kernel B: 16 KiB of registers and 16 KiB of shared memory per
% block.  Each filled segment is one resident block's allocation; white
% is idle.  Scale 0.45 mm per KiB.
% Original (no outline figure for this lesson).
% Usage: % Split versus unified on-chip storage (Lesson 6, own example).  An SM
% with a 128 KiB register file and 64 KiB of shared memory, or one
% 192 KiB pool.  Kernel A: 64 KiB of registers per block, no shared
% memory.  Kernel B: 16 KiB of registers and 16 KiB of shared memory per
% block.  Each filled segment is one resident block's allocation; white
% is idle.  Scale 0.45 mm per KiB.
% Original (no outline figure for this lesson).
% Usage: \input{figures/l06-unified-pool}
% Width: body (120 mm); the drawing is about 116 mm wide.
\begin{tikzpicture}[x=1mm, y=1mm, font=\scriptsize\sffamily]
  \def\k{0.45}
  % \seg{x0 KiB}{y}{size KiB}{fill}
  \newcommand{\seg}[4]{\draw[fill=#4] ({26+(#1)*\k},#2) rectangle ++({(#3)*\k},4);}
  \newcommand{\idle}[3]{\draw[ln-rule, fill=white] ({26+(#1)*\k},#2) rectangle ++({(#3)*\k},4);}
  % headings
  \node[anchor=south west, text=ln-muted] at (26,46) {\iflangja{レジスタファイル 128\,KiB}{register file 128 KiB}};
  \node[anchor=south west, text=ln-muted] at (26+131*\k,46) {\iflangja{共有メモリ 64\,KiB}{shared 64 KiB}};
  % Kernel A split: 2 blocks x 64 KiB regs
  \node[anchor=west] at (0,42) {\iflangja{A：分離}{A, split}};
  \seg{0}{40}{64}{ln-accent!30} \seg{64}{40}{64}{ln-accent!30}
  \idle{131}{40}{64}
  \node[anchor=west] at (26+196*\k,42) {2};
  % Kernel A unified
  \node[anchor=west] at (0,34) {\iflangja{A：統合}{A, unified}};
  \seg{0}{32}{64}{ln-accent!30} \seg{64}{32}{64}{ln-accent!30} \seg{128}{32}{64}{ln-accent!30}
  \node[anchor=west] at (26+196*\k,34) {3};
  % Kernel B split: 4 blocks x 16 regs + 4 x 16 smem
  \node[anchor=west] at (0,22) {\iflangja{B：分離}{B, split}};
  \foreach \b in {0,...,3} { \seg{\b*16}{20}{16}{ln-accent!30} \seg{131+\b*16}{20}{16}{ln-tip!35} }
  \idle{64}{20}{64}
  \node[anchor=west] at (26+196*\k,22) {4};
  % Kernel B unified: 6 blocks x (16 + 16)
  \node[anchor=west] at (0,14) {\iflangja{B：統合}{B, unified}};
  \foreach \b in {0,...,5} { \seg{\b*32}{12}{16}{ln-accent!30} \seg{\b*32+16}{12}{16}{ln-tip!35} }
  \node[anchor=west] at (26+196*\k,14) {6};
  \node[anchor=south, text=ln-muted] at (26+198*\k,46) {\iflangja{ブロック}{blocks}};
  % legend
  \draw[fill=ln-accent!30] (26,3) rectangle ++(4,3);
  \node[anchor=west] at (31,4.5) {\iflangja{レジスタ}{registers}};
  \draw[fill=ln-tip!35] (50,3) rectangle ++(4,3);
  \node[anchor=west] at (55,4.5) {\iflangja{共有メモリ}{shared memory}};
  \draw[ln-rule, fill=white] (78,3) rectangle ++(4,3);
  \node[anchor=west] at (83,4.5) {\iflangja{未使用}{idle}};
  \node[anchor=south west, text=ln-muted] at (26,26.5)
    {\iflangja{統合プール 192\,KiB（1 本の帯）}{unified pool 192 KiB (one bar)}};
\end{tikzpicture}

% Width: body (120 mm); the drawing is about 116 mm wide.
\begin{tikzpicture}[x=1mm, y=1mm, font=\scriptsize\sffamily]
  \def\k{0.45}
  % \seg{x0 KiB}{y}{size KiB}{fill}
  \newcommand{\seg}[4]{\draw[fill=#4] ({26+(#1)*\k},#2) rectangle ++({(#3)*\k},4);}
  \newcommand{\idle}[3]{\draw[ln-rule, fill=white] ({26+(#1)*\k},#2) rectangle ++({(#3)*\k},4);}
  % headings
  \node[anchor=south west, text=ln-muted] at (26,46) {\iflangja{レジスタファイル 128\,KiB}{register file 128 KiB}};
  \node[anchor=south west, text=ln-muted] at (26+131*\k,46) {\iflangja{共有メモリ 64\,KiB}{shared 64 KiB}};
  % Kernel A split: 2 blocks x 64 KiB regs
  \node[anchor=west] at (0,42) {\iflangja{A：分離}{A, split}};
  \seg{0}{40}{64}{ln-accent!30} \seg{64}{40}{64}{ln-accent!30}
  \idle{131}{40}{64}
  \node[anchor=west] at (26+196*\k,42) {2};
  % Kernel A unified
  \node[anchor=west] at (0,34) {\iflangja{A：統合}{A, unified}};
  \seg{0}{32}{64}{ln-accent!30} \seg{64}{32}{64}{ln-accent!30} \seg{128}{32}{64}{ln-accent!30}
  \node[anchor=west] at (26+196*\k,34) {3};
  % Kernel B split: 4 blocks x 16 regs + 4 x 16 smem
  \node[anchor=west] at (0,22) {\iflangja{B：分離}{B, split}};
  \foreach \b in {0,...,3} { \seg{\b*16}{20}{16}{ln-accent!30} \seg{131+\b*16}{20}{16}{ln-tip!35} }
  \idle{64}{20}{64}
  \node[anchor=west] at (26+196*\k,22) {4};
  % Kernel B unified: 6 blocks x (16 + 16)
  \node[anchor=west] at (0,14) {\iflangja{B：統合}{B, unified}};
  \foreach \b in {0,...,5} { \seg{\b*32}{12}{16}{ln-accent!30} \seg{\b*32+16}{12}{16}{ln-tip!35} }
  \node[anchor=west] at (26+196*\k,14) {6};
  \node[anchor=south, text=ln-muted] at (26+198*\k,46) {\iflangja{ブロック}{blocks}};
  % legend
  \draw[fill=ln-accent!30] (26,3) rectangle ++(4,3);
  \node[anchor=west] at (31,4.5) {\iflangja{レジスタ}{registers}};
  \draw[fill=ln-tip!35] (50,3) rectangle ++(4,3);
  \node[anchor=west] at (55,4.5) {\iflangja{共有メモリ}{shared memory}};
  \draw[ln-rule, fill=white] (78,3) rectangle ++(4,3);
  \node[anchor=west] at (83,4.5) {\iflangja{未使用}{idle}};
  \node[anchor=south west, text=ln-muted] at (26,26.5)
    {\iflangja{統合プール 192\,KiB（1 本の帯）}{unified pool 192 KiB (one bar)}};
\end{tikzpicture}

% Width: body (120 mm); the drawing is about 116 mm wide.
\begin{tikzpicture}[x=1mm, y=1mm, font=\scriptsize\sffamily]
  \def\k{0.45}
  % \seg{x0 KiB}{y}{size KiB}{fill}
  \newcommand{\seg}[4]{\draw[fill=#4] ({26+(#1)*\k},#2) rectangle ++({(#3)*\k},4);}
  \newcommand{\idle}[3]{\draw[ln-rule, fill=white] ({26+(#1)*\k},#2) rectangle ++({(#3)*\k},4);}
  % headings
  \node[anchor=south west, text=ln-muted] at (26,46) {\iflangja{レジスタファイル 128\,KiB}{register file 128 KiB}};
  \node[anchor=south west, text=ln-muted] at (26+131*\k,46) {\iflangja{共有メモリ 64\,KiB}{shared 64 KiB}};
  % Kernel A split: 2 blocks x 64 KiB regs
  \node[anchor=west] at (0,42) {\iflangja{A：分離}{A, split}};
  \seg{0}{40}{64}{ln-accent!30} \seg{64}{40}{64}{ln-accent!30}
  \idle{131}{40}{64}
  \node[anchor=west] at (26+196*\k,42) {2};
  % Kernel A unified
  \node[anchor=west] at (0,34) {\iflangja{A：統合}{A, unified}};
  \seg{0}{32}{64}{ln-accent!30} \seg{64}{32}{64}{ln-accent!30} \seg{128}{32}{64}{ln-accent!30}
  \node[anchor=west] at (26+196*\k,34) {3};
  % Kernel B split: 4 blocks x 16 regs + 4 x 16 smem
  \node[anchor=west] at (0,22) {\iflangja{B：分離}{B, split}};
  \foreach \b in {0,...,3} { \seg{\b*16}{20}{16}{ln-accent!30} \seg{131+\b*16}{20}{16}{ln-tip!35} }
  \idle{64}{20}{64}
  \node[anchor=west] at (26+196*\k,22) {4};
  % Kernel B unified: 6 blocks x (16 + 16)
  \node[anchor=west] at (0,14) {\iflangja{B：統合}{B, unified}};
  \foreach \b in {0,...,5} { \seg{\b*32}{12}{16}{ln-accent!30} \seg{\b*32+16}{12}{16}{ln-tip!35} }
  \node[anchor=west] at (26+196*\k,14) {6};
  \node[anchor=south, text=ln-muted] at (26+198*\k,46) {\iflangja{ブロック}{blocks}};
  % legend
  \draw[fill=ln-accent!30] (26,3) rectangle ++(4,3);
  \node[anchor=west] at (31,4.5) {\iflangja{レジスタ}{registers}};
  \draw[fill=ln-tip!35] (50,3) rectangle ++(4,3);
  \node[anchor=west] at (55,4.5) {\iflangja{共有メモリ}{shared memory}};
  \draw[ln-rule, fill=white] (78,3) rectangle ++(4,3);
  \node[anchor=west] at (83,4.5) {\iflangja{未使用}{idle}};
  \node[anchor=south west, text=ln-muted] at (26,26.5)
    {\iflangja{統合プール 192\,KiB（1 本の帯）}{unified pool 192 KiB (one bar)}};
\end{tikzpicture}

  \caption{ブロックあたり256スレッドの2つのカーネルにおける分離型と
    統合型の記憶．カーネル~A はレジスタファイルを埋め，共有メモリを
    遊ばせる．カーネル~B は共有メモリを埋め，レジスタファイルの半分を
    遊ばせる．1つの \SI{192}{\kibi\byte} のプールなら A は3ブロック，
    B は6ブロックを受け入れる．}
  \label{fig:l06-unified-pool}
\end{figure}

もし \SI{192}{\kibi\byte} が1つのプールであり，各ブロックがそこから
レジスタと共有メモリを取るなら，カーネル~A は $192/64 = 3$ ブロック，
カーネル~B は $192/32 = 6$ ブロックが収まる．どちらも常駐スレッドが
50\,\% 増える（\cref{tab:l06-unified}）．統合設計では，占有率はレジスタと
共有メモリによって\emph{合わせて}制限される．

\begin{margintable}
  \caption{\cref{fig:l06-unified-pool}のカーネルの SM あたり常駐
    ブロック数（スレッド数）．}
  \label{tab:l06-unified}
  \begin{tabular}{@{}lcc@{}}
    \toprule
    カーネル & 分離 & 統合 \\
    \midrule
    A & 2 (512) & 3 (768) \\
    B & 4 (1024) & 6 (1536) \\
    \bottomrule
  \end{tabular}
\end{margintable}

\subsection{キャッシュ，スクラッチパッド，レジスタファイル}

レジスタと共有メモリが記憶を共有できるかどうかは，それぞれがどう
アドレス指定されるかによる．キャッシュとスクラッチパッドは，何を保持するかを
誰が決めるかで異なる．

\begin{definition}[ハードウェア管理キャッシュとソフトウェア管理キャッシュ]
  \label{def:l06-managed-caches}
  \term[はーどうぇあかんりきゃっしゅ]{ハードウェア管理キャッシュ}[hardware-managed cache]
  は，どのメモリブロックを保持するかを自ら決める．したがって各エントリは
  \term[たぐ]{タグ}[tag]，すなわち保持するブロックのアドレスの上位ビットを
  持ち，要求されたアドレスがあるかどうかを知るためにアクセスのたびに
  比較される．満杯のキャッシュは
  \term[ちかんぽりし]{置換ポリシ}[replacement policy]が選んだブロックを
  追い出す．たとえば
  \term[えるあーるゆー]{LRU}[LRU (least recently used)]\index{least recently used|see{LRU}}は，最後の
  アクセスが最も過去にあるブロックを追い出す．
  \term[そふとうぇあかんりきゃっしゅ]{ソフトウェア管理キャッシュ}[software-managed cache]
  （第3章のスクラッチパッドメモリ，CUDA の共有メモリ）はプログラムが
  置いたものを保持する．プログラムはすべてのデータがどこにあるかを
  知っているので，タグもミスも置換もない．
\end{definition}
\margin{NVIDIA GPU はグローバルメモリを，\SI{32}{\byte} のセクタ4つから
成る \SI{128}{\byte} のラインでキャッシュし（第4章），ワープのアクセスは
\SI{32}{\byte} のトランザクションで処理される（BPG §10.2.1 Coalesced
Access to Global Memory \cite{nvidia2026cudabestpractices}）．}

レジスタファイルはスクラッチパッドに似ている．どちらもデータだけを格納し，
タグを持たず，メモリアドレスではなく位置によってエントリを選ぶ．
両者は2つの点で異なる（\cref{tab:l06-storage}）．\margin{共有メモリの
レイテンシは，キャッシュされないグローバルメモリのおよそ100分の1である
\cite{harris2013sharedmem}．}算術命令は2つのソースレジスタ（融合積和演算
$d = a \cdot b + c$ では3つ）を読み，1つに書き込むので，1回のアクセスで
1つのアドレスを処理するメモリよりも多くのポートが必要である．また，
命令が使うレジスタは命令の中で名指しされるので，そのアクセスパターンは
実行前に静的に分かる．一方スクラッチパッドのアドレスは実行時に計算される．
この静的な知識こそが，帯域が高いにもかかわらず，第4章のバンク配置と
上述のコンパイラ管理の ORF を可能にする．

\begin{table}[tb]
  \centering
  \caption{3種類のオンチップ記憶．}
  \label{tab:l06-storage}
  \setlength{\tabcolsep}{4pt}
  \begin{tabular}{@{}>{\raggedright\arraybackslash}p{34mm}
                  >{\raggedright\arraybackslash}p{24mm}
                  >{\raggedright\arraybackslash}p{26mm}
                  >{\raggedright\arraybackslash}p{28mm}@{}}
    \toprule
    & ハードウェアキャッシュ（L1） & スクラッチパッド（共有メモリ） & レジスタファイル \\
    \midrule
    内容を決めるもの & ハードウェア & プログラム & コンパイラ \\
    タグ & あり & なし & なし \\
    位置の特定 & アドレスとタグ & 位置 & 位置（レジスタ名） \\
    \addlinespace
    命令あたりのアクセス & 1アドレス & 1アドレス & 2--3読み出し，1書き込み \\
    パターンが分かる時点 & 実行時 & 実行時 & コンパイル時 \\
    \bottomrule
  \end{tabular}
\end{table}

\subsection{統合された構造}

3つの構造はすべて同じ SM 上の SRAM アレイなので，1つのアレイで
それらすべてを受け持ち，カーネルごとに分割できる．\margin{Gebhart らは，
\enquote{reduce accesses to the main register file by 60\,\%}という
以前のレジスタ階層と2レベルスケジューラの上に構築しており，そのため
専用のレジスタファイルよりバイトあたりのポートが少ないプールでも
オペランドを供給できる \cite{gebhart2012unifying}．}

\begin{definition}[統合ローカルメモリ]
  \label{def:l06-unified-local-memory}
  \term[とうごうろーかるめもり]{統合ローカルメモリ}[unified local memory]
  と呼ばれる設計は，SM ごとの単一のオンチップ SRAM プールであり，
  その容量はカーネルごとにレジスタファイル，共有メモリ，そして完全な提案では
  \term[えるわんでーたきゃっしゅ]{L1データキャッシュ}[L1 data cache]，
  すなわちグローバルメモリとローカルメモリへのアクセスのための
  ハードウェア管理の1次キャッシュに分けられる．
\end{definition}

統合プールは各利用者の最も厳しい要求を満たさなければならない．すなわち
レジスタファイルの帯域，したがって十分な読み出しポートと書き込みポート，
そしてどちらの構造単独よりも多くのバンクである．レジスタアクセス
（全レーンで同じレジスタ）と共有メモリアクセス（計算されたアドレスによる）
はパターンが異なり，どちらも空いたバンクを見つけなければならないからである．
その代わりに，各カーネルは必要な分割を得て，占有率の制限が緩和される．

\caution{講義 vs 論文：本章は講義に従い，レジスタファイルと共有メモリを
統合する．Gebhart らは\emph{3つの}構造，\enquote{the register file,
scratchpad (shared memory in NVIDIA terminology), and cache} を統合し，
\enquote{a performance benefit as high as 71\,\% along with an energy
reduction up to 33\,\%}を報告している \cite{gebhart2012unifying}．}%
出荷されている GPU が1つのアレイを共有させているのは L1 データキャッシュと
共有メモリの間だけであり（Fermi がそうであり，Volta 以降の GPU が再び
そうしている），レジスタファイルは別のままである．

\begin{definition}[共有メモリのカーブアウト]
  \label{def:l06-carveout}
  \term[きょうゆうめもりのかーぶあうと]{共有メモリのカーブアウト}[shared memory carve-out]
  とは，SM の L1 データキャッシュと共有メモリを合わせたアレイのうち，
  共有メモリとして構成される部分である．残りは L1 キャッシュとして働く．
  CUDA ランタイムはそれを，カーネルの共有メモリの使用量と任意の
  優先指定に基づいて，固定されたサイズの集合からカーネルごとに選ぶ．
\end{definition}
\tip{ブロックあたり \SI{48}{\kibi\byte} を超える動的共有メモリを必要とする
カーネルは，\texttt{cudaFuncAttribute\allowbreak MaxDynamic\allowbreak SharedMemorySize}
で明示的に有効にしなければならない（CPG，Compute Capability 7.x，Shared
Memory \cite{nvidia2024cudaguide}）．}%
優先指定は1回のランタイム呼び出しでカーネルごとに設定する．

\Needspace{6\baselineskip}% ja only: keeps the short listing on one page
\begin{lstlisting}[language=C++, numbers=none]
// カーネル k に共有メモリの最大容量の約半分を
// 希望する（ヒントであり，ランタイムは対応する
// カーブアウトに切り上げる）．
cudaFuncSetAttribute(k,
    cudaFuncAttributePreferredSharedMemoryCarveout, 50);
\end{lstlisting}

\margin{SM あたりの L1 と共有メモリの合計：V100 で \SI{128}{\kibi\byte}，
A100 で \SI{192}{\kibi\byte}，H100 で \SI{256}{\kibi\byte}．カーブアウト：
V100 は 0，8，16，32，64，\SI{96}{\kibi\byte}，A100 は最大
\SI{164}{\kibi\byte}，H100 は 0，8，16，32，64，100，132，164，196，
\SI{228}{\kibi\byte}（CPG，Compute Capability 7.x，8.x，9.0
\cite{nvidia2024cudaguide}）．}%
こうして共有メモリをほとんど使わないカーネルは大きな L1 を得，データを
タイル化するカーネルはアレイの大部分をスクラッチパッドとして得る．
これは\cref{fig:l06-unified-pool}の柔軟性を2つの構造について実現したもの
である．

\begin{keypoint}[章のまとめ]
  ダイバージェント分岐では，ワープは経路を部分的なマスクで順に実行し，
  直接後支配節点で再収束する．if変換は分岐を，両方の経路をフェッチする
  プレディケート付きのコードに変える．SIMT スタックは，入れ子のものも
  含めダイバージェント分岐で（代替 PC，マスク，再収束 PC）のエントリを
  プッシュし，再収束 PC でポップする．幅の広いワープはフェッチを節約するが
  ダイバージェンスが増える．動的ワープ形成は1つの PC にある異なるワープの
  スレッドを，それぞれホームレーンに置いたままマージする．非生存の値，
  短命の値，一度だけ読まれる値が，レジスタファイルキャッシュ，ORF/LRF の
  階層，高速部と低速部への分割を正当化する．仮想化は名前をより少ない
  物理レジスタに対応付ける．統合ローカルメモリは1つのプールをカーネル
  ごとにレジスタ，共有メモリ，L1 に分割する．現在の GPU は L1 と共有メモリに
  ついて，カーブアウトによってそれを行っている．
\end{keypoint}

\end{document}
